\documentclass[11pt]{amsart}
\usepackage{lmodern}
\usepackage{amsmath,amssymb,amsthm,mathtools}
\usepackage[hmargin=1.1in,vmargin=1in]{geometry}
\usepackage{xcolor}
\usepackage{booktabs}
\usepackage{enumitem}
\usepackage[colorlinks=true,linkcolor=blue!50!black,citecolor=blue!50!black,urlcolor=blue!50!black]{hyperref}
\hypersetup{pdftitle={Weighted pair energy and the sixth moment of the Riemann zeta function},pdfauthor={Min Chae}}


\newtheorem{theorem}{Theorem}[section]
\newtheorem{proposition}[theorem]{Proposition}
\newtheorem{lemma}[theorem]{Lemma}
\newtheorem{corollary}[theorem]{Corollary}
\theoremstyle{definition}
\newtheorem{definition}[theorem]{Definition}
\theoremstyle{remark}
\newtheorem{remark}[theorem]{Remark}

\numberwithin{equation}{section}

\newcommand{\cN}{\mathcal{N}}
\newcommand{\cJ}{\mathcal{J}}
\newcommand{\cK}{\mathcal{K}}
\newcommand{\Mst}{M_{*}}
\DeclareMathOperator{\sgn}{sgn}

\title[Weighted pair energy and the sixth moment]{Weighted pair energy and the sixth moment of the Riemann zeta function}
\author{Min Chae}
\date{15 November 2026; revised 17 November 2026}

\begin{document}
\raggedbottom
\setlength{\emergencystretch}{2em}
\renewcommand{\arraystretch}{1.35}

\begin{abstract}
We prove
\[
\int_0^T|\zeta(\tfrac12+it)|^6\,dt\ll_\varepsilon T^{84599/68640+\varepsilon},
\qquad \frac{84599}{68640}=1.2325029\ldots .
\]
We use the Atkinson method of de Faveri and Pandey and the curved
model of bramble\_trace. In the reciprocal kernel, we keep the weight
$D^{k-1/2}$ and sum it using the refined pair count. This gives the
sixth-moment bound with Heath-Brown's exponent pair
$(1/245,\,269/280+\delta)$.
Starting from the curved exponential sum, we derive the curved energy
estimate, exact phase, and signed reciprocal kernel of that model.
We also prove the Gaussian form of Atkinson's formula.
We use the sparse fourth-moment theorem and the optimal Vinogradov
mean-value theorem as inputs.
\end{abstract}

\maketitle

%======================================================================
\section{Introduction}
%======================================================================

\subsection{Background}
For $T\ge2$ let
\[
M_6(T)=\int_0^T|\zeta(\tfrac12+it)|^6\,dt .
\]
The conjectured order of $M_6(T)$ is $T(\log T)^9$ \cite{CG98}.
The classical bound $M_6(T)\ll_\varepsilon T^{5/4+\varepsilon}$ follows
from H\"older's inequality, Ingham's fourth moment
$M_4(T)\ll T(\log T)^4$ \cite{Ing26}, and Heath-Brown's twelfth moment
$M_{12}(T)\ll T^2(\log T)^{17}$ \cite{HB78}.
De Faveri and Pandey \cite{dFP} improved $5/4$ to $37/30$.
They used Atkinson's formula, frequency spacing, and close pairs of
large values. Paper~\cite{B38} gives the exponent
$12867/10438=1.23270\ldots$ using a curved Atkinson phase and a
reciprocal kernel. Sections~\ref{sec:curved}--\ref{sec:kernel}
give proofs of the parts of that argument used here.

\subsection{Results}
For $T\ge2$ and $V>0$ let $\cN(T,V)$ be the largest cardinality of a $1$-separated set of $t\in[T,2T]$ with $|\zeta(\frac12+it)|\ge V$. We write $\mathcal X\prec\mathcal Y$ if $\mathcal X\ll_\varepsilon T^{\varepsilon}\mathcal Y$ for every $\varepsilon>0$, the implied constant depending only on $\varepsilon$ and on fixed parameters.

\begin{theorem}[sixth moment]\label{thm:sixth}
For $T\ge2$ and every $\varepsilon>0$,
\[
\int_0^T|\zeta(\tfrac12+it)|^6\,dt\ll_\varepsilon T^{84599/68640+\varepsilon}.
\]
\end{theorem}

The exponent improves $37/30$ by $19/22880$ and $12867/10438$ by $73259/358232160$, that is, by $0.000204\ldots$. Theorem~\ref{thm:sixth} follows from the large-value estimate below and the fourth and twelfth moments.

\begin{theorem}[large values]\label{thm:large}
Let $\sigma_0=15959/137280$. Uniformly for $\sigma_0\le\sigma\le1/8$,
\[
\cN(T,T^\sigma)\prec T^{(54928-293156\sigma)/38969}.
\]
\end{theorem}

Theorem~\ref{thm:large} is the case $(k,\ell)=(1/245,269/280)$ of a general statement. Fix an exponent pair $(k,\ell)$ and put $a=1-\ell+k$. For $N,W\ge1$ define the two length costs
\begin{equation}\label{eq:lammu}
\lambda=\Big(\frac TN\Big)^{k}\Big(\frac WN\Big)^{-a},\qquad
\mu=W^{k}\Big(\frac WN\Big)^{-a},
\end{equation}
and the nine quantities
\begin{equation}\label{eq:nine}
\begin{aligned}
B_1&=T^{1/2}N^{3/2}W^{-1}V^{-2}, &
B_2&=\lambda^{1/4}T^{1/4}N^{5/4}W^{-1/4}V^{-2},\\
B_3&=\mu^{1/2}T^{-1/4}N^{5/4}W^{1/2}V^{-2}, &
B_4&=\mu^{1/2}T^{-1/2}N^{3/2}W^{3/4}V^{-2},\\
B_5&=\big(\lambda^{4}T^{6+4k}N^{20}W^{-4}V^{-44-24k}\big)^{1/(17+2k)}, &
B_6&=\lambda^{1/4}T^{-1/8}N^{11/8}W^{1/4}V^{-2},\\
B_7&=T^{1/3}N^{2/3}V^{-2}, &
B_8&=T^{1/2}N^{25/48}W^{1/6}V^{-10/3},\\
B_9&=(TN/V^{6})^{4/7}.
\end{aligned}
\end{equation}

\begin{theorem}[large values, general exponent pair]\label{thm:general}
Let $0\le k\le\frac12\le\ell\le1$ be such that $(k,\ell+\delta)$ is an exponent pair in the sense of Definition~\ref{def:pair} for every sufficiently small $\delta>0$, and let $a=1-\ell+k$ satisfy
\begin{equation}\label{eq:H}
0<k<\tfrac12,\qquad a<\tfrac12+k,\qquad 8a<3\big(k+\tfrac12\big).
\end{equation}
Let $\frac1{10}<\sigma\le\frac18$, $\nu_*=1-4\sigma$, and let $q\in[\frac98,\frac43]$ satisfy
\begin{equation}\label{eq:P}
\big(q-\tfrac12\big)\nu_*<\tfrac12,\quad
(3-q)\nu_*<1,\quad
(2+q)\nu_*<2,\quad
\big(q+\tfrac12\big)\nu_*-\tfrac12<4\sigma .
\end{equation}
Then, with $V=T^\sigma$,
\[
\cN(T,V)\prec\max_{T^{1/3}\le N\le T^{\nu_*}}\ \sum_{j=1}^{9}B_j\big|_{W=N^{q}} .
\]
The implied constants depend on $(k,\ell)$, on $\varepsilon$ and on the margins in \eqref{eq:P}.
\end{theorem}

\begin{remark}\label{rem:eighth}
For $E(\sigma)=(54928-293156\sigma)/38969$, we have
$\frac{d}{d\sigma}\big(E(\sigma)+8\sigma\big)=18596/38969>0$.
The eighth-moment layer therefore reaches its maximum at the right
endpoint. The eighth-moment paper~\cite{Thin} treats that range using
thin planar energy and an upper-amplitude bound.
\end{remark}

\subsection{Relation to previous work}\label{sec:new}
We use the Atkinson reduction, frequency spacing, clustering, and
refined pair count of \cite{dFP}. The curved model and signed reciprocal
kernel were introduced in \cite{B38}. We prove each part used here from
the sum \eqref{eq:model}. The derivations cover the curvature
energy, exact critical phase, Morse--Fresnel formula, both Sobolev
classes, and the signed sum over Poisson modes. They are given in
Sections~\ref{sec:curved}--\ref{sec:kernel} and
Appendices~\ref{app:energyproof}--\ref{app:dualityproof}.
The factor $F^k$ and the stationary amplitude $H^{-1/2}$ follow from
these derivations.

We make three changes.
\begin{enumerate}[label=(\roman*)]
\item We keep two bounds for the same kernel simultaneously (Corollary~\ref{cor:two}), one of which keeps the weight $D^{k-1/2}$ in the additive discrepancy $D$. This leads to the weighted pair energy
\[
A_k(S)=\sup_{1\le D\ll U}D^{k-1/2}E_2(S;D)
\]
in the large-value estimate for the curved model (Proposition~\ref{prop:model}).
\item We bound $A_k$ using the refined pair count of \cite{dFP}.
The bounds $E_2\ll R_1D$ and $E_2\prec R_1(1+A_g)$ meet at
$D\asymp1+A_g$. We sum the weight $D^k$ at this scale
(Lemma~\ref{lem:weighted}).
\item We keep both length factors $\lambda,\mu$ in \eqref{eq:lammu}
through the zeta reduction (Theorem~\ref{thm:general}).
We then choose Heath-Brown's pair $j=8$.
\end{enumerate}
These steps give the large-value bound in Theorem~\ref{thm:general}.

\subsection{Proof outline}\label{sec:status}
Appendices~\ref{app:elementary}--\ref{app:dualityproof} prove the
spacing, curved energy, Morse--Fresnel formula, and signed block bound.
Appendix~\ref{app:transferproof} proves the clustering and pair-count
steps. Appendix~\ref{app:upperproof} gives the upper-amplitude bound.

The sixth-moment proof uses two input theorems: the sparse fourth
moment~\cite{Iwa80} and the optimal Vinogradov mean
value~\cite{BDG16,Woo16}. Their statements are
given in Appendix~\ref{app:base}. Appendix~\ref{app:atkinsondirect}
proves the short-second-moment estimate by Gaussian smoothing.
Appendix~\ref{app:localzeta} proves the local zeta inequality and
derives the fourth and twelfth moments from the sparse fourth moment.
Appendix~\ref{app:hbderivative} derives the derivative estimate from
the Vinogradov mean value. The chosen pair uses derivatives through
order 12 and mean values through degree 11.

%======================================================================
%======================================================================
\section{Notation and inputs}
%======================================================================

\subsection{Notation}
$e(x)=e^{2\pi ix}$. A set $X\subset\mathbb R$ is $1$-separated if distinct points differ by at least $1$. For a finite multiset $X\subset\mathbb R$ and $D>0$ put
\begin{align*}
E_2(X;D)&=\#\{(x,y)\in X^2:|x-y|\le D\},\\
E_4(X;D)&=\#\{(x_1,\dots,x_4)\in X^4:|x_1+x_2-x_3-x_4|\le D\},
\end{align*}
counted with multiplicity. Thus $E_4(X;D)=E_2(X\!+\!X;D)$, where $X\!+\!X$ is the multiset of sums $x_1+x_2$. A sum over $D$ is over $D=2^j\ge1$.

\begin{lemma}\label{lem:enlarge}
Let $X$ be a finite multiset and $D>0$. Then $\#\{(x,y):|x-y-s|\le D/2\}\le 2E_2(X;D)$ for every $s\in\mathbb R$, and $E_2(X;cD)\le(4c+6)E_2(X;D)$ for $c\ge1$. If $X$ is $1$-separated, $E_2(X;D)\le(2D+1)|X|$.
\end{lemma}

\begin{proof}
Let $\psi(x)=\max(0,1-|x|)$, so $\psi\ge\frac12$ on $[-\frac12,\frac12]$, $0\le\psi\le\mathbf 1_{[-1,1]}$ and $\widehat\psi(\xi)=(\sin\pi\xi/\pi\xi)^2\ge0$. With $\mathcal F(\xi)=\sum_{x\in X}e(x\xi)$,
\begin{align*}
\#\{|x-y-s|\le D/2\}&\le2\sum_{x,y}\psi\Big(\frac{x-y-s}{D}\Big)
=2D\int\widehat\psi(D\xi)e(-s\xi)|\mathcal F(\xi)|^2d\xi\\
&\le2\sum_{x,y}\psi\Big(\frac{x-y}{D}\Big)\le2E_2(X;D).
\end{align*}
Covering $[-cD,cD]$ by $2c+3$ intervals of length $D$ gives the second bound. The third is immediate.
\end{proof}

\subsection{Exponent pairs}

\begin{definition}\label{def:pair}
We use the finite-derivative convention of \cite[\S1.1]{TY}.
Let $0\le k\le\frac12\le\ell\le1$. We say $(k,\ell)$ is an exponent pair if for every $s>0$ there are $R\in\mathbb N$ and $\epsilon_0>0$ such that the following holds. Let $M\ge1$, $I\subset(M,2M]$ an interval, $y\ge M^s$, and let $g\in C^{R+1}(I)$ satisfy
\[
\Big|g^{(r+1)}(x)-\frac{d^r}{dx^r}\big(yx^{-s}\big)\Big|\le\epsilon_0\Big|\frac{d^r}{dx^r}\big(yx^{-s}\big)\Big|
\qquad(x\in I,\ 0\le r\le R).
\]
Then $\sum_{m\in I}e(g(m))\ll (y/M^{s})^{k}M^{\ell}$.
\end{definition}

\begin{lemma}[Heath-Brown {\cite[Theorem~2]{HB17}}]\label{lem:hb}
For every integer $j\ge3$ and every sufficiently small $\delta>0$,
\[
\Big(\frac2{(j-1)^2(j+2)},\ 1-\frac{3j-2}{j(j-1)(j+2)}+\delta\Big)
\]
is an exponent pair, with $R$ bounded on every range in which $\log(y/M^{s})/\log M$ is bounded. For $j=8$ this is $(1/245,269/280+\delta)$, and $a=1-\ell+k=17/392$.
\end{lemma}
For the pair $j=8$ used below, Appendix~\ref{app:base} proves the
finite-derivative conditions from the derivative estimate
\eqref{eq:basederivative}, using orders at most 12.

\subsection{Moments}
The classical moments \cite{Ing26,HB78}, together with the local zeta estimate \eqref{eq:baselocal}, give $\cN(T,V)\prec TV^{-4}$ and $\cN(T,V)\prec T^2V^{-12}$. Appendix~\ref{app:localzeta} derives these bounds from the sparse fourth-moment theorem and the local zeta inequality.

\subsection{The reduction of de Faveri and Pandey}

\begin{proposition}[clustering and refined pair count; {\cite[\S\S7--9]{dFP}}]\label{in:dfp}
Let $\frac1{10}<\sigma\le\frac18$, $V=T^\sigma$, and let $W\ge1$ be a block length with $W\le\sqrt{TN}$ for the $N$ below. There are $F_c\ge1$, a scale $N$ with
\[
T^{1/3}\ll N\prec T/V^4,
\]
$K\prec T/(V^2W)$ blocks with initial heights $t_i\asymp T$, coefficients $|a^{(i)}_n|\le1$ of the form described in Lemma~\ref{lem:atkinson}, and $1$-separated sets $S_i\subset[0,W]$ with $|D_{h_i}(w)|\ge Z$ on $S_i$, such that, with $R_1=\sum_i|S_i|$,
\begin{equation}\label{eq:cluster}
Z\succ F_c(TN)^{1/4},\qquad \cN(T,V)\prec ZR_1(TN)^{-1/4},
\end{equation}
and for all $D\ge1$
\begin{equation}\label{eq:paircount}
\sum_iE_2(S_i;D)\prec R_1+\frac{R_1D\sqrt{T/N}}{F_cV^4}+\frac{T\sqrt{R_1}}{F_c^{3/2}V^6}.
\end{equation}
\end{proposition}
Proposition~\ref{prop:clusterproof} proves this reduction from the
theorems stated in Appendix~\ref{app:base}; the range is restricted to
$\sigma>1/10$ to justify the first phase truncation.

For the upper-amplitude range we will use the estimate proved in Appendix~\ref{app:upperproof}:
\begin{equation}\label{eq:dfpline}
\cN(T,T^\sigma)\prec T^{(94-500\sigma)/67}
\qquad\left(\frac18\le\sigma\le\frac5{38}\right).
\end{equation}
Appendix~\ref{app:upperproof} derives this estimate from the model and the stated zeta inputs.

%======================================================================
\section{The curved model and its curvature energy}\label{sec:curved}
%======================================================================

Let $N\ge2$, $W\ge1$, $\eta,\kappa\in\mathbb R$, $|a_n|\le1$, and
\begin{equation}\label{eq:model}
D_h(w)=\sum_{N<n\le2N}a_n\,e\Big(w\Big(\frac nN\Big)^{1/2}+h(w)\Big(\frac nN\Big)^{3/2}\Big),
\qquad h(w)=\frac{\eta w}{1+\kappa w}.
\end{equation}
In this section and the next, $\prec$ allows losses $N^{\varepsilon}$. Throughout we assume
\begin{equation}\label{eq:modelranges}
N^{9/8}\le W\le N^{4/3},\qquad 0<|\eta|\le c\,W/N^2,\qquad |\eta|N^{1/8}+|\kappa|W\le c,
\end{equation}
with $c>0$ small enough in terms of the fixed parameters. Let $S\subset[0,W]$ be $1$-separated, $r=|S|$, and $|D_h(w)|\ge Z\ge1$ on $S$. Put
\[
U=1+|\eta|W,\qquad B=|\eta\kappa|W^2,\qquad X=N^{1/8},\qquad H_0=C\sqrt{WX},\qquad \Mst=W/N .
\]

\begin{lemma}[curvature multiplier; {\cite[\S2]{B38}}]\label{lem:mult}
Let $\xi_n=(n/N)^{1/2}+\eta(n/N)^{3/2}$, $\chi_n=(n/N)^{3/2}$, $r_0(w)=-\eta\kappa w^2/(1+\kappa w)$, and let $\beta$ be a fixed smooth bump. Then
\[
K(u):=\sum_{m\ne n}\Big|\widehat{\beta(\cdot/W)\,e\big((\chi_m-\chi_n)r_0(\cdot)\big)}\big(u-(\xi_m-\xi_n)\big)\Big|\prec N^2\sqrt{1+B|u|}.
\]
\end{lemma}

\begin{proof}
See Proposition~\ref{prop:multproof}, including the independent spacing proof in Lemma~\ref{lem:spacingproof}.
\end{proof}

\begin{proposition}[curvature energy; {\cite[Theorem~2]{B38}}]\label{prop:curv}
For $1\le D\ll U$,
\begin{equation}\label{eq:curv}
E_4(S;D)\prec\Big(\frac NZ\Big)^4\Big(1+\frac BD\Big)E_2(S;D)+\frac N{Z^2}\,Dr^3 .
\end{equation}
\end{proposition}

\begin{proof}
This is the case $j=2$ of Proposition~\ref{prop:allenergy}, whose proof starts directly with the large-value condition and keeps all fixed coefficient classes.
\end{proof}

%======================================================================
\section{Duality, Poisson summation and the exact tuple phase}\label{sec:signed}
%======================================================================

\subsection{The signed tuple expansion}
Fix a dyadic $M$ with $1\le M\le\Mst$, a pair of neighboring height blocks, and one translation of the heights, chosen independently of the Poisson index $m$. The translated model is again of the form \eqref{eq:model}, with $\eta,\kappa$ replaced by $\eta(1+\kappa u)^{-2}$ and $\kappa(1+\kappa u)^{-1}$; these are uniformly comparable with $\eta,\kappa$. Let $b_w$ ($|b_w|=1$) be the duality coefficients; they depend on $S$ and on the coefficients $a_n$ and remain fixed across $m$. In the Sobolev step both classes $f_{m,0}$ and $f_{m,1}=\partial_zf_{m,0}$ are kept, where
\[
f_{m,0}(z)=\sum_wb_w\,a_{m,0}(z,w)\,e(\Theta_m(z,w)).
\]
After the fourth powers are opened, the quantity to be bounded is, for each sign of $m$ and each class,
\begin{equation}\label{eq:tuple}
\cK_M:=\sum_{\mathbf w\in S^4}b_{w_1}b_{w_2}\overline{b_{w_3}b_{w_4}}\;K_M(\mathbf w),\qquad
K_M(\mathbf w)=\sum_{|m|\asymp M}\int A_{\mathbf w}(m/M,x)\,e\big(\Phi(m/M,x)\big)\,dx,
\end{equation}
where $\Phi(t,x)=\sum_{i=1}^4\epsilon_i\Theta_{Mt}(NMtx,w_i)$, $\epsilon=(1,1,-1,-1)$. The absolute value is taken after the inner signed sum over $m$, with the same coefficients $b_w$ for every $m$.

\begin{proposition}[duality and Poisson, with explicit blocks]\label{in:duality}
For each fixed neighbor, sign and Sobolev class, let $\mathfrak K_M$ be the sum over the height blocks of the absolute values of their signed tuple kernels, as defined in Corollary~\ref{cor:blockconsumer}. Then
\begin{equation}\label{eq:duality}
r^2Z^2\prec N^2r+r^{3/4}\sum_{M,h,\pm,\ell}
N^{7/4}M^{1/2}\mathfrak K_M^{1/4}.
\end{equation}
\end{proposition}
The term $N^2r$ bounds the zero and nonstationary modes. The stationary part has an exact representation.
\begin{proof}
Proposition~\ref{prop:dualityproof} derives the block inequality from the sum \eqref{eq:model}. Corollary~\ref{cor:blockconsumer} applies H\"older to the block sum. Each block uses its own translated $Q$.
\end{proof}

\subsection{The exact critical phase}
Put $y=Nm$, $\Delta=(z-w)/y$, $\eta_{\rm eff}=\eta/((1+\kappa z)(1+\kappa w))$ and $\tau=\eta_{\rm eff}\Delta^2$.

\begin{lemma}[{\cite[(3.10)--(3.12)]{B38}}]\label{lem:phase}

The critical value of the Poisson phase is
\[
\Theta_m(z,w)=\frac{y\Delta^2}{12}\,\Omega(\tau)\big(4-\Omega(\tau)\big),\qquad
\Omega(\tau)=\frac2{1+\sqrt{1-3\tau}} .
\]
In particular
\[
\Theta_m(z,w)=\frac{(z-w)^2}{4Nm}+\frac{\eta(z-w)^4}{8(Nm)^3(1+\kappa z)(1+\kappa w)}+O\big(y|\tau|^2\Delta^2\big).
\]
\end{lemma}
\begin{proof}
Proposition~\ref{prop:exactphaseproof} solves the critical quadratic, identifies its admissible branch, and gives the exact analytic remainder with all needed derivatives.
\end{proof}

\begin{remark}
The expansion is consistent with the closed form: $\Omega(\tau)=1+\frac34\tau+O(\tau^2)$, hence $\Omega(4-\Omega)=3+\frac32\tau+O(\tau^2)$. We keep the full remainder in the phase.
\end{remark}

\begin{lemma}[two-pair differences]\label{lem:twopair}
Let $f$ be real analytic on a neighborhood of a compact interval $I$. For $v_1,\dots,v_4\in I$,
\[
|f(v_1)+f(v_2)-f(v_3)-f(v_4)|\ll_f|v_1+v_2-v_3-v_4|+|v_1^2+v_2^2-v_3^2-v_4^2|,
\]
and the same holds, uniformly, after any fixed number of derivatives in parameters on which $f$ depends analytically.
\end{lemma}

\begin{proof}
Write $v_1=s/2+\delta$, $v_2=s/2-\delta$. The function $\delta\mapsto f(s/2+\delta)+f(s/2-\delta)$ is even and analytic, so $f(v_1)+f(v_2)=G(s,d)$ with $d=(v_1-v_2)^2$ and $G$ analytic on a neighborhood of the relevant compact set. Let $s'=v_3+v_4$, $d'=(v_3-v_4)^2$. Then $s-s'=P:=v_1+v_2-v_3-v_4$ and
\[
d-d'=2Q-(s-s')(s+s'),\qquad Q:=v_1^2+v_2^2-v_3^2-v_4^2,
\]
so $|d-d'|\ll|Q|+|P|$. The mean value theorem for $G$ gives the bound; derivatives in parameters are handled identically.
\end{proof}

\subsection{The tuple phase}
Let $L=NM$, $t=m/M$,
\[
P=w_1+w_2-w_3-w_4,\qquad Q=w_1^2+w_2^2-w_3^2-w_4^2,\qquad F=|Q|/L,\qquad D=1+|P|.
\]
On the stationary support all $w_i$ lie in an interval of length $O(L)$ after the translation, so $F\ll NM\le W$. By Lemmas~\ref{lem:phase} and \ref{lem:twopair}, with $z=Nmx$,
\begin{equation}\label{eq:psi}
\begin{gathered}
\Phi(t,x)=\frac{Q}{4Lt}+\psi(t,x),\qquad
\psi_x=-\frac P2+O\big(|\eta|(F+|P|)\big),\\
\psi_{xx}=\frac{3\eta Q}{2Nm}-3\eta Px+O\big((|\kappa|L+|\eta|)|\eta|(F+|P|)\big),
\end{gathered}
\end{equation}
with the same control for every fixed number of further derivatives. The explicit mixed-derivative inequalities and their proof are in Proposition~\ref{prop:exactphaseproof}.

\begin{lemma}[Morse--Fresnel; {\cite[\S3]{B38}}]\label{lem:morse}

Let $H=|\eta|F$. If $H\ge1$ and $|P|\ll H$, then
\[
\int A_{\mathbf w}(t,x)\,e(\Phi(t,x))\,dx=H^{-1/2}b(t)\,e(S(t)),\qquad
S(t)=\frac Q{4Lt}+R(t),
\]
with $R^{(j)}(t)\ll_j|\eta|F$ and $b^{(j)}(t)\ll_j1$ for every fixed $j$. Here $H$ is fixed as $t$ varies.
\end{lemma}
\begin{proof}
See Lemma~\ref{lem:exactmorse} and Proposition~\ref{prop:tuplemorseproof}.
\end{proof}

If $H\le1$ the correction $\psi$ can be absorbed into the amplitude. If $|P|\gg1+H$, integration by parts in $x$, repeated a fixed number of times more than the number of $t$-derivatives needed, gives an amplitude decaying like $(D/(1+H))^{-J}$. In both cases the phase is the exact reciprocal $Q/(4Lt)$.

%======================================================================
\section{The kernel estimate}\label{sec:kernel}
%======================================================================

\begin{proposition}\label{prop:kernel}
Let $(k,\ell)$ satisfy the hypotheses of Theorem~\ref{thm:general}, let $\delta>0$ be fixed, and apply Definition~\ref{def:pair} to the pair $(k,\ell+\delta)$. For $M\ge N^{1/24}$, every fixed $J$, each sign of $m$ and each Sobolev class,
\begin{equation}\label{eq:kernel}
|K_M(\mathbf w)|\ll M^{1-a+\delta}F^{k}(1+|\eta|F)^{-1/2}\Big(1+\frac D{1+|\eta|F}\Big)^{-J}\qquad(F\ge X),
\end{equation}
and $|K_M(\mathbf w)|\ll_J MD^{-J}$ for $F<X$.
\end{proposition}

\begin{proof}
\emph{Case $F\ge4M$.} By Lemma~\ref{lem:morse} and the remarks after it, the inner integral equals $(1+H)^{-1/2}$ times an amplitude with bounded $t$-derivatives and decay $(1+D/(1+H))^{-J}$, times $e(S(m/M))$, where $S$ is the exact reciprocal phase plus a correction whose $j$-th derivative is a relative $O_j(|\eta|)$ of that of $Q/(4Lt)$. With $g(m)=-\sgn(Q)S(m/M)$ the first derivative is modeled by $ym^{-2}$, $y=|Q|/(4N)=FM/4$, and $y/M^2=F/(4M)\ge1$. For the fixed exponent pair, Definition~\ref{def:pair} gives a finite derivative order $R$ at $s=2$. The relative errors $O_j(|\eta|)$ for $j\le R+1$ are below its $\epsilon_0$ once the model smallness constant is chosen accordingly. Definition~\ref{def:pair} gives $\sum_{m\in I}e(g(m))\ll(F/M)^kM^{\ell+\delta}=M^{1-a+\delta}F^k$ for every subinterval $I$. Partial summation restores the amplitude. For $j=8$ only derivatives through order 12 are needed, as shown in Appendix~\ref{app:base}.

\emph{Case $X\le F\le4M$.} Here $H\le4|\eta|M\ll W^2/N^3\le N^{-1/3}$, so the phase is the exact reciprocal and the amplitude is smooth. The second-derivative test gives $\ll\sqrt F+M/\sqrt F$. Since $F\le4M$ and $a<\frac12+k$, $\sqrt F\ll F^{k}M^{1/2-k}\le F^kM^{1-a}$. Since $M\le\Mst\le N^{1/3}$, $F\ge N^{1/8}$ and $8a<3(k+\frac12)$, $M^{a}\le N^{a/3}\le N^{(k+1/2)/8}\le F^{k+1/2}$, i.e.\ $M/\sqrt F\le M^{1-a}F^k$.

\emph{Case $F<X$.} Then $H\le|\eta|X=o(1)$, the phase in $x$ has derivative $-P/2+o(D)$, and integration by parts gives $D^{-J}$; summing trivially over $m$ gives $MD^{-J}$.
\end{proof}

\begin{corollary}[two simultaneous bounds]\label{cor:two}
For $M\ge N^{1/24}$, $F\ge X$ and every fixed $A$, both
\begin{align*}
|K_M(\mathbf w)|&\ll M^{1-a+\delta}(NM)^{k}D^{-1/2}\Big(1+\frac DU\Big)^{-A}
\qquad\text{and}\\
|K_M(\mathbf w)|&\ll M^{1-a+\delta}|\eta|^{-k}D^{k-1/2}\Big(1+\frac DU\Big)^{-A}
\end{align*}
hold for the same tuple $\mathbf w$.
\end{corollary}

\begin{proof}
Let $H=|\eta|F$. Since $F\ll NM\le W$, $1+H\ll U$. If $D\le1+H$, then $(1+H)^{-1/2}\le D^{-1/2}$ and $F^k\ll(NM)^k$ give the first bound, while $F^k(1+H)^{-1/2}=|\eta|^{-k}H^{k}(1+H)^{-1/2}\le|\eta|^{-k}(1+H)^{k-1/2}\le|\eta|^{-k}D^{k-1/2}$ (as $k<\frac12$) gives the second; in this range $1+D/U\asymp1$. If $D>1+H$, then
\[
F^k(1+H)^{-1/2}\Big(\frac{1+H}D\Big)^{J}\le\min\Big((NM)^kD^{-1/2},\,|\eta|^{-k}D^{k-1/2}\Big)\Big(\frac{1+H}D\Big)^{J-1/2},
\]
and $(1+H)/D\ll(1+D/U)^{-1}$; take $J\ge A+\frac12$.
\end{proof}

The two bounds are properties of the same kernel; we use each on a different part of the energy (Proposition~\ref{prop:model}). The dependence between $F$ and $P$ is kept.

\begin{lemma}[small $F$]\label{lem:smallF}
For every fixed $J>3$,
\[
\sum_{\substack{\mathbf w\in S^4\\ F<X}}D^{-J}\ll rE_2(S;C\sqrt{LX}),
\]
uniformly in all parameters.
\end{lemma}

\begin{proof}
Let $b=w_1-w_3$, $c=w_1-w_4$. Then $w_2=P+w_1-b-c$ and
\[
Q=2bc+2P(w_1-b-c)+P^2 .
\]
If $|P|\le D'$ and $|Q|\le LX$, then $|bc|\ll L(X+D')$ (recall $|w_i|\ll L$ and $D'\ll L$), so $\min(|b|,|c|)\ll\sqrt{L(X+D')}$. Choose the close pair (at most $2E_2(S;C\sqrt{L(X+D')})$ ways), the remaining point among $w_3,w_4$ ($r$ ways), and $w_2$ in an interval of length $O(D')$ ($O(1+D')$ ways). By Lemma~\ref{lem:enlarge}, $E_2(S;C\sqrt{L(X+D')})\ll(1+\sqrt{D'})E_2(S;C\sqrt{LX})$. Summing over dyadic $D'$ against $D'^{-J}$ gives the claim.
\end{proof}

%======================================================================
\section{The large-value estimate for the curved model}\label{sec:model}
%======================================================================

Let $(k,\ell)$ satisfy \eqref{eq:H}, $a=1-\ell+k$, and define, for the model,
\[
\lambda_0=|\eta|^{-k}\Mst^{-a},\qquad \mu_0=N^{k}\Mst^{k-a}=W^{k}\Mst^{-a},\qquad
A_k(S)=\sup_{1\le D\ll U}D^{k-1/2}E_2(S;D).
\]

\begin{proposition}[strengthened curved large values]\label{prop:model}
Assume \eqref{eq:modelranges} and the exponent-pair hypothesis. Then
\begin{equation}\label{eq:modelbound}
\begin{aligned}
r\prec{}&\frac{N^2}{Z^2}
+\lambda_0^{1/5}\frac{N^{8/5}W^{3/5}}{Z^{12/5}}A_k(S)^{1/5}
+\mu_0^{1/2}\frac{N^{5/2}W^{3/2}}{Z^5}U^{1/4}\\
&+\lambda_0^{1/4}B^{1/4}\frac{N^2W^{3/4}}{Z^3}
+\frac{NW^{3/4}}{Z^2}E_2(S;H_0)^{1/4}.
\end{aligned}
\end{equation}
\end{proposition}

\begin{proof}
Put
\[
\cJ=\max_{M,h,\pm,\ell}\Big(\frac M{\Mst}\Big)^{3}M^{-1}\mathfrak K_M,
\]
with the block sum $\mathfrak K_M$ of Corollary~\ref{cor:blockconsumer}. The following energy sums are first taken on translated blocks and then bounded by their global counterparts, as proved in that corollary. Since
\[
N^{7/4}M^{1/2}\big(\Mst^3M^{-2}\cJ\big)^{1/4}=N^{7/4}\Mst^{3/4}\cJ^{1/4}=NW^{3/4}\cJ^{1/4},
\]
Proposition~\ref{in:duality} gives
\begin{equation}\label{eq:rJ}
r^2Z^2\prec N^2r+NW^{3/4}r^{3/4}\cJ^{1/4}.
\end{equation}
We bound $\cJ$. Fix $M$.

\emph{Tuples with $F\ge X$, $M\ge N^{1/24}$.} Split the tuples by $D=1+|P|$ into dyadic layers. A layer is contained in $\{|P|\le 2D\}$, of size at most $E_4(S;2D)\ll E_4(S;D)$ by Lemma~\ref{lem:enlarge}. For $D\gg U$ the factor $(1+D/U)^{-A}$ and $E_4(S;D)\ll(D/U)E_4(S;U)$ reduce to $D\asymp U$. For $D\ll U$ insert \eqref{eq:curv}. On the parts $(N/Z)^4E_2(S;D)$ and $(N/Z)^4(B/D)E_2(S;D)$ use the second bound of Corollary~\ref{cor:two}; on $(N/Z^2)Dr^3$ use the first. This choice is legitimate because both bounds hold for every tuple. With Lemma~\ref{lem:enlarge} ($E_2(S;D)\ll rD$) and $k<\frac12$,
\[
\sum_D D^{k-1/2}E_2(S;D)\prec A_k(S),\qquad
B\sum_DD^{k-3/2}E_2(S;D)\ll Br,\qquad
\sum_{D\ll U}D^{-1/2}Dr^3\ll r^3\sqrt U .
\]
Hence these tuples contribute to $\sum_{\mathbf w}|K_M|$ at most
\[
M^{1-a+\delta}\Big[|\eta|^{-k}\Big(\frac NZ\Big)^4\big(A_k(S)+Br\big)+(NM)^k\frac N{Z^2}r^3\sqrt U\Big].
\]
Multiplying by $(M/\Mst)^3M^{-1}$ and using $3-a>0$, $3-a+k>0$ and $M\le\Mst$:
\[
\Big(\frac M\Mst\Big)^3M^{-a}|\eta|^{-k}\le\lambda_0,\qquad
\Big(\frac M\Mst\Big)^3M^{-a}(NM)^k\le N^k\Mst^{k-a}=\mu_0 .
\]
(The loss $M^\delta$ is absorbed in $\prec$ after fixing $\delta$ small.)

\emph{Tuples with $F<X$.} By Proposition~\ref{prop:kernel} and Lemma~\ref{lem:smallF}, with $L=NM\le W$, their contribution to $(M/\Mst)^3M^{-1}\sum|K_M|$ is $\ll rE_2(S;C\sqrt{WX})=rE_2(S;H_0)$.

\emph{$M<N^{1/24}$.} Proposition~\ref{prop:kernel} with the trivial bound in $m$ in place of the exponent pair (i.e.\ the case $k=0$, $a=0$ of the argument of Corollary~\ref{cor:two}) gives $|K_M|\ll MD^{-1/2}(1+D/U)^{-A}$. The weight is $(M/\Mst)^3\le N^{1/8}\Mst^{-3}=:C_M$. Since $\Mst\ge N^{1/8}$, $|\eta|\le1$ and $a<2$,
\[
C_M/\lambda_0\le N^{1/8}\Mst^{a-3}\le N^{(a-2)/8}<1,\qquad
C_M/\mu_0=N^{1/8-k}\Mst^{a-k-3}\le N^{(a-2-9k)/8}<1,
\]
and $D^{-1/2}\le D^{k-1/2}$, so these $M$ are covered by the previous bound.

Altogether
\begin{equation}\label{eq:Jbound}
\cJ\prec\Big(\frac NZ\Big)^4\lambda_0A_k(S)+\Big(\frac NZ\Big)^4\lambda_0Br+\frac N{Z^2}\mu_0r^3\sqrt U+rE_2(S;H_0).
\end{equation}
Insert \eqref{eq:Jbound} into \eqref{eq:rJ}; one of the following holds.
\begin{itemize}
\item $r^2Z^2\prec N^2r$: $r\prec N^2Z^{-2}$.
\item $r^2Z^2\prec NW^{3/4}r^{3/4}(N/Z)\lambda_0^{1/4}A_k^{1/4}$: $r^{5/4}\prec\lambda_0^{1/4}N^2W^{3/4}Z^{-3}A_k^{1/4}$, the second term.
\item $r^2Z^2\prec NW^{3/4}r^{3/4}(N^{1/4}Z^{-1/2})\mu_0^{1/4}r^{3/4}U^{1/8}$: $r^{1/2}\prec\mu_0^{1/4}N^{5/4}W^{3/4}Z^{-5/2}U^{1/8}$, the third term.
\item $r^2Z^2\prec NW^{3/4}r^{3/4}(N/Z)\lambda_0^{1/4}B^{1/4}r^{1/4}$: the fourth term.
\item $r^2Z^2\prec NW^{3/4}r^{3/4}r^{1/4}E_2(S;H_0)^{1/4}$: the fifth term.\qedhere
\end{itemize}
\end{proof}

%======================================================================
\section{Return to the zeta function}\label{sec:return}
%======================================================================

\subsection{Atkinson's phase and the curved model}

\begin{lemma}\label{lem:atkinson}
Use the normalization of \cite[(41)]{dFP},
\[
f(t,n)=\frac t\pi\left(\operatorname{arsinh}\sqrt{\frac{\pi n}{2t}}+
\sqrt{\frac{\pi n}{2t}\left(1+\frac{\pi n}{2t}\right)}\right).
\]
The constant phase $-1/8$ in the usual cosine form is absorbed into the coefficients. For $n\asymp N\le T^{3/5}$ and $t\asymp T$,
\[
f(t,n)=\sqrt{\frac{2nt}\pi}+\frac{\sqrt{2\pi}}{12}\frac{n^{3/2}}{\sqrt t}+r_t(n),\qquad
|r_t(n)|+N|\partial_nr_t(n)|\ll N^{5/2}T^{-3/2}.
\]
For a block starting at $t_0\asymp T$ put
\[
w=\sqrt{2N/\pi}\,\big(\sqrt t-\sqrt{t_0}\big),\qquad \eta=-\frac{\pi N}{12t_0},\qquad \kappa=\Big(\frac{2Nt_0}\pi\Big)^{-1/2}.
\]
Then the first two terms of $f(t,n)-f(t_0,n)$ equal exactly $w(n/N)^{1/2}+h(w)(n/N)^{3/2}$ with $h(w)=\eta w/(1+\kappa w)$. In particular $|\eta|\asymp N/T$, $|\kappa|\asymp(TN)^{-1/2}$, $U\asymp1+NW/T$ and $B\asymp N^{1/2}W^2T^{-3/2}$.
\end{lemma}

\begin{proof}

With $x=\sqrt{\pi n/(2t)}$ the phase of Atkinson's formula with the constant phase removed is $2t\operatorname{arsinh}x+\sqrt{2\pi nt+\pi^2n^2}$. Expanding, its $\frac1{2\pi}$-multiple is $\sqrt{2nt/\pi}+\frac{\sqrt{2\pi}}{12}n^{3/2}t^{-1/2}+O(n^{5/2}t^{-3/2})$, with the stated derivative bound. Since $\sqrt t=\sqrt{t_0}+w\sqrt{\pi/(2N)}$, the first term contributes $\sqrt{2n/\pi}\,(\sqrt t-\sqrt{t_0})=w(n/N)^{1/2}$. For the second,
\[
\frac{\sqrt{2\pi}}{12}n^{3/2}\Big(\frac1{\sqrt t}-\frac1{\sqrt{t_0}}\Big)
=-\frac{\sqrt{2\pi}}{12}n^{3/2}\frac{w\sqrt{\pi/(2N)}}{\sqrt{t_0}\big(\sqrt{t_0}+w\sqrt{\pi/(2N)}\big)}
=-\frac{\pi N}{12t_0}\,\frac{w}{1+\kappa w}\Big(\frac nN\Big)^{3/2}.
\]
The last statements follow from $B=|\eta\kappa|W^2$.
\end{proof}

Lemma~\ref{lem:completionproof} proves the partial-summation and interval-completion steps, with a common Fourier multiplier in $n$. Proposition~\ref{prop:clusterproof} tracks the divisor loss, the selected level set and all spacing classes.

\subsection{The weighted pair energy}

\begin{lemma}\label{lem:weighted}
In the setting of Proposition~\ref{in:dfp}, assume $\sqrt{T/N}+W\sqrt{N/T}\prec V^4$ and $0<k<\frac12$. Put $A_g=T/(F_c^{3/2}V^6\sqrt{R_1})$. Then
\[
\sum_iA_k(S_i)\prec R_1(1+A_g)^{k+1/2}.
\]
\end{lemma}

\begin{proof}
For $D\ll U\asymp1+NW/T$, the middle term of \eqref{eq:paircount} is at most $R_1U\sqrt{T/N}/V^4\prec R_1$ by the assumption. Hence, with Lemma~\ref{lem:enlarge},
\[
\sum_iE_2(S_i;D)\prec R_1\min(D,1+A_g)\qquad(1\le D\ll U).
\]
Bounding the supremum by the sum over dyadic $D$,
\begin{align*}
\sum_iA_k(S_i)&\le\sum_{D\ll U}D^{k-1/2}\sum_iE_2(S_i;2D)\\
&\prec R_1\Big(\sum_{D\le1+A_g}D^{k+1/2}+(1+A_g)\sum_{D>1+A_g}D^{k-1/2}\Big)\ll R_1(1+A_g)^{k+1/2},
\end{align*}
since $k+\frac12>0>k-\frac12$.
\end{proof}

Lemma~\ref{lem:weighted} sums the weight $D^k$ at $D\asymp1+A_g$, where the two pair-count bounds meet.

\subsection{Proof of Theorem~\ref{thm:general}}
Let $N$ be as in Proposition~\ref{in:dfp} and $W=N^q$. For $T^{1/3}\le N\le T^{\nu_*}$ the hypotheses \eqref{eq:P} give \eqref{eq:modelranges} with $c\to0$ as $T\to\infty$ (see the dictionary in Table~\ref{tab:cond}), $W\le\sqrt{TN}$, $N\le T^{3/5}$ (as $\sigma>\frac1{10}$), and the assumption of Lemma~\ref{lem:weighted}: $\sqrt{T/N}\le T^{1/3}<V^4$ since $\sigma>\frac1{12}$, and $W\sqrt{N/T}\prec V^4$ by the last condition in \eqref{eq:P}. Each of these quantities is $T^{(\text{affine in }\nu)}$ with positive $\nu$-coefficient, so it suffices to check them at $\nu=\nu_*$. By Lemma~\ref{lem:atkinson}, $\lambda_0\asymp\lambda$ and $\mu_0=\mu$ as in \eqref{eq:lammu}.

Apply Proposition~\ref{prop:model} to each $S_i$ and sum over $i$. We use $K\prec T/(V^2W)$, H\"older's inequality over $i$, \eqref{eq:cluster}, Lemma~\ref{lem:weighted} and \eqref{eq:paircount}. In each case below we bound $R_1$, then $\cN(T,V)\prec ZR_1(TN)^{-1/4}$; the resulting power of $Z$ is negative, so we may insert $Z\succ F_c(TN)^{1/4}$, and the resulting power of $F_c$ is nonpositive, so $F_c\ge1$ may be dropped at the end.

\emph{Term 1.} $R_1\prec KN^2Z^{-2}$, so $\cN\prec KN^{7/4}T^{-1/4}Z^{-1}\prec KN^{3/2}T^{-1/2}\prec B_1$.

\emph{Terms 2 and 5.} By H\"older, $\sum_ir_i\le C^{4/5}\sum_iA_k(S_i)^{1/5}\le C^{4/5}K^{4/5}(\sum_iA_k(S_i))^{1/5}$ with $C=\lambda^{1/4}N^2W^{3/4}Z^{-3}$. By Lemma~\ref{lem:weighted},
\[
R_1\prec K\lambda^{1/4}N^2W^{3/4}Z^{-3}(1+A_g)^{(1+2k)/8}.
\]
The constant part of $(1+A_g)$ gives $\cN\prec K\lambda^{1/4}N^{7/4}W^{3/4}T^{-1/4}Z^{-2}\prec B_2$. The part $A_g^{(1+2k)/8}$ gives
\[
R_1^{(17+2k)/16}\prec K\lambda^{1/4}N^2W^{3/4}Z^{-3}\Big(\frac T{F_c^{3/2}V^6}\Big)^{(1+2k)/8}.
\]
Then $\cN\prec ZR_1(TN)^{-1/4}$ carries $Z^{(2k-31)/(17+2k)}$; inserting $Z\succ(TN)^{1/4}$ and $K\prec T/(V^2W)$ and raising to the power $17+2k$ gives
\[
\cN^{17+2k}\prec(TN)^{-12}\,K^{16}\lambda^4N^{32}W^{12}\Big(\frac T{V^6}\Big)^{2+4k}=\lambda^4T^{6+4k}N^{20}W^{-4}V^{-44-24k},
\]
that is, $\cN\prec B_5$. Here $\lambda$ keeps its full factor $(T/N)^k$ together with its power of $W/N$.

\emph{Terms 3 and 4.} $R_1\prec K\mu^{1/2}N^{5/2}W^{3/2}Z^{-5}U^{1/4}$, so $\cN\prec K\mu^{1/2}U^{1/4}N^{9/4}W^{3/2}T^{-1/4}Z^{-4}\prec K\mu^{1/2}U^{1/4}N^{5/4}W^{3/2}T^{-5/4}$. With $U^{1/4}\ll1+(NW/T)^{1/4}$ this is $\ll B_3+B_4$.

\emph{Term 6.} $R_1\prec K\lambda^{1/4}B^{1/4}N^2W^{3/4}Z^{-3}$, so $\cN\prec K\lambda^{1/4}B^{1/4}N^{5/4}W^{3/4}T^{-3/4}$; with $B^{1/4}\asymp N^{1/8}W^{1/2}T^{-3/8}$ this is $B_6$.

\emph{Terms 7, 8, 9.} By H\"older, $R_1\le NW^{3/4}Z^{-2}K^{3/4}(\sum_iE_2(S_i;H_0))^{1/4}$, and \eqref{eq:paircount} at $D=H_0\asymp W^{1/2}N^{1/16}$ has three terms. The term $R_1$ gives $R_1\prec KN^{4/3}WZ^{-8/3}$ and $\cN\prec KN^{2/3}WT^{-2/3}=B_7$. The term $R_1H_0\sqrt{T/N}/(F_cV^4)$ gives
\[
\cN\prec KN^{2/3}WT^{-2/3}\big(H_0\sqrt{T/N}\,V^{-4}\big)^{1/3}=B_8 .
\]
The term $T\sqrt{R_1}/(F_c^{3/2}V^6)$ gives $R_1^{7/8}\prec NW^{3/4}Z^{-2}K^{3/4}T^{1/4}V^{-3/2}$, whence
\[
\cN\prec(TN)^{-4/7}\big(NW^{3/4}K^{3/4}T^{1/4}V^{-3/2}\big)^{8/7}=B_9 .
\]

The exponents of $F_c$ in the nine terms are $-1,-2,-4,-4,-2,-2,-5/3,-2,-12/7$. This proves Theorem~\ref{thm:general}. \qed

%======================================================================
\section{Optimization: proof of Theorem~\ref{thm:large}}\label{sec:optimize}
%======================================================================

Take $(k,\ell)=(1/245,269/280)$ (Lemma~\ref{lem:hb} with $j=8$), so $a=17/392$. The hypotheses \eqref{eq:H} hold: $a<\frac12+k$ and $8a=17/49<3(k+\frac12)=1.512\ldots$. Write $N=T^\nu$, $W=T^\omega$, $V=T^\sigma$. Then
\[
\log_T\lambda=k+(a-k)\nu-a\omega,\qquad \log_T\mu=k\omega-a(\omega-\nu),
\]
and $\log_TB_j=e_j(\sigma,\nu,\omega)$ is affine. Put $\nu_*=1-4\sigma$ and
\[
\omega(\sigma)=\frac{5(1692769+2461308\sigma)}{15938321},\qquad q(\sigma)=\frac{\omega(\sigma)}{\nu_*},\qquad W=N^{q(\sigma)} .
\]
This choice solves $e_4=e_5$ at $\nu=\nu_*$, and there
\[
e_4=e_5=E(\sigma):=\frac{54928-293156\sigma}{38969}.
\]
The function $q$ is increasing, with $q(\sigma_0)=8714085/7509649$ and $q(1/8)=20004325/15938321$, so $\frac98<q<\frac43$ on $[\sigma_0,\frac18]$. For fixed $q$ in this range every $e_j$ has positive $\nu$-coefficient, so the maximum over $N$ in Theorem~\ref{thm:general} is attained at $\nu=\nu_*$. At $\nu=\nu_*$ each $E-e_j$ is affine in $\sigma$, and Table~\ref{tab:margins} gives its values at both endpoints.

\begin{table}[ht]
\centering
\caption{Margins $E(\sigma)-e_j$ at $\nu=\nu_*$.}\label{tab:margins}
\begin{tabular}{ccc}
\toprule
$j$ & $\sigma=\sigma_0$ & $\sigma=1/8$\\
\midrule
1 & $301109/3509220$ & $81147/838859$\\
2 & $22819/5174400$ & $15613/8039920$\\
3 & $1093427/28073760$ & $1016501/31876642$\\
4 & $0$ & $0$\\
5 & $0$ & $0$\\
6 & $29384059/13756142400$ & $10383337/15619554580$\\
7 & $15959/205920$ & $24557/467628$\\
8 & $27213367/673770240$ & $31884151/1530078816$\\
9 & $4519/80080$ & $3165/77938$\\
\bottomrule
\end{tabular}
\end{table}

The conditions \eqref{eq:P} and \eqref{eq:modelranges} are verified in Table~\ref{tab:cond}; each quantity is increasing in $\nu$ for fixed $q$ and affine in $\sigma$ at $\nu=\nu_*$.

\begin{table}[ht]
\centering
\caption{Parameter conditions: $\log_T$ of each quantity at $\nu=\nu_*$.}\label{tab:cond}
\begin{tabular}{lccc}
\toprule
quantity & $\log_T$ & $\sigma=\sigma_0$ & $\sigma=1/8$\\
\midrule
$|\kappa|W$ & $(q-\frac12)\nu-\frac12$ & $-4118359/28073760$ & $-7806313/63753284$\\
$|\eta|N^2/W$ & $(3-q)\nu-1$ & $-37003/2339480$ & $-2033002/15938321$\\
$|\eta|N^{1/8}$ & $\frac98\nu-1$ & $-36437/91520$ & $-7/16$\\
$N^{5/2}T^{-3/2}$ & $\frac52\nu-\frac32$ & $-2231/13728$ & $-1/4$\\
$\eta^2W$ & $(2+q)\nu-2$ & $-4340377/14036880$ & $-11872317/31876642$\\
$W\sqrt{N/T}\,V^{-4}$ & $(q+\frac12)\nu-\frac12-4\sigma$ & $-717841/9357920$ & $-7806313/63753284$\\
\bottomrule
\end{tabular}
\end{table}

The phase formula is exact. The bounds $|\eta|\le cW/N^2$ and $W\le N^{4/3}$ also give $\eta^2W\le c^2$.

Since every margin is uniform on $[\sigma_0,\frac18]$, the constants are uniform in $\sigma$. Theorem~\ref{thm:general} gives $\cN(T,T^\sigma)\prec T^{E(\sigma)}$, which is Theorem~\ref{thm:large}. The intersection with the fourth-moment line is $E(\sigma_0)=1-4\sigma_0$, $\sigma_0=15959/137280$. \qed

\begin{remark}
On the whole range the active terms are $B_4$ (the growth $U^{1/4}$ of the diagonal) and $B_5$ (the refined pair count). A further improvement needs either control of how many pairs of large values occupy the radius $D\asymp1+A_g$, or additional cancellation in the tuple sum that keeps the product of the four dual coefficients.
\end{remark}

%======================================================================
\section{Proof of Theorem~\ref{thm:sixth}}\label{sec:moment}
%======================================================================

Decompose $[T,2T]$ according to the dyadic size $V=T^\sigma$ of $|\zeta(\frac12+it)|$. The unit intervals meeting $\{t:|\zeta(\frac12+it)|\ge V\}$ split into two classes (even and odd), and choosing one point in each interval of a class gives a $1$-separated set; hence the measure of this set is at most $2\cN(T,V)$. The contribution of the layer $V$ to $\int_T^{2T}|\zeta|^6$ is therefore $\prec V^6\cN(T,V)=T^{6\sigma}\cN(T,T^\sigma)$, and we bound the exponent of this layer.
\begin{itemize}
\item $\sigma\le\sigma_0$: by the fourth moment, $1+2\sigma\le1+2\sigma_0=84599/68640$.
\item $\sigma_0\le\sigma\le\frac18$: by Theorem~\ref{thm:large}, $E(\sigma)+6\sigma$, which decreases (slope $6-293156/38969<0$) from the value $84599/68640$ at $\sigma_0$.
\item $\frac18\le\sigma\le\frac5{38}$: by Proposition~\ref{prop:curvatureline}, $(94-500\sigma)/67+6\sigma\le327/268<84599/68640$.
\item $\sigma\ge\frac5{38}$: by the twelfth moment, $2-6\sigma\le23/19<84599/68640$.
\end{itemize}
Amplitudes $|\zeta|\le1$ and bounded $t$ contribute $O(T)$. Summing over the $O(\log T)$ layers and over dyadic $T$ proves Theorem~\ref{thm:sixth}. The order of choices is: $\varepsilon$; the exponent-pair loss $\delta$; the numbers of derivatives and integrations by parts; $T_0$. The fixed positive margins of the inactive terms and the fixed negative margins in Table~\ref{tab:cond} absorb all logarithmic and $T^{o(1)}$ losses. \qed

%======================================================================
\appendix
\section{Elementary estimates and frequency spacing}\label{app:elementary}

We prove the sampling, derivative, and spacing estimates used below.
The spacing estimate follows the argument of de Faveri and Pandey.
Appendix~\ref{app:base} states the zeta and mean-value theorems used in
the moment bounds.

\begin{lemma}[local sampling]\label{lem:localproof}
If $p\ge1$, $f\in C^1$, and $I_x=[x-1/3,x+1/3]$, then
\[
 |f(x)|^p\ll_p\int_{I_x}(|f(u)|^p+|f'(u)|^p)\,du.
\]
Hence the same inequality can be summed over any $1$-separated set,
with a constant uniform in its cardinality.
\end{lemma}
\begin{proof}
For every $u\in I_x$, the fundamental theorem of calculus gives
$|f(x)|\le |f(u)|+\int_{I_x}|f'(v)|\,dv$. Raise this inequality to the
$p$th power, use H\"older on the integral, and average over $u$.
The intervals $I_x$ for distinct sample points have disjoint interiors.
\end{proof}

\begin{lemma}[second derivative estimates]\label{lem:secondproof}
On an interval of length at most $M$, suppose
$0<\lambda\le |g''|\le C\lambda$. Then
\[
 \left|\sum_{n\in I}e(g(n))\right|
 \ll_C M\sqrt\lambda+\lambda^{-1/2}+1.
\]
For an integral with a compactly supported $C^1$ amplitude $a$, if
$|g''|\ge\lambda>0$, then
\[
 \left|\int a(x)e(g(x))\,dx\right|
 \ll\lambda^{-1/2}(\|a\|_\infty+\|a'\|_1).
\]
\end{lemma}
\begin{proof}
The integral estimate follows by removing $|g'|\le\sqrt\lambda$, whose
length is at most $2/\sqrt\lambda$, and integrating by parts on each of
the two remaining intervals. Since $g'$ is monotone,
$\int |g''|/|g'|^2\le2/\sqrt\lambda$ there.

For the sum, the statement is trivial for $\lambda\ge1/16$. Otherwise
put $\delta=\sqrt\lambda$. The range of the monotone function $g'$
meets $O_C(M\lambda+1)$ integers. Remove the inverse images of the
$2\delta$-neighborhoods of these integers and the unit intervals meeting
their endpoints. They contain
$O_C((M\lambda+1)(\delta/\lambda+1))$ integers. On each complementary
interval the increments $g(n+1)-g(n)$ are monotone and lie between
$q+\delta$ and $q+1-\delta$ for one integer $q$. The identity
\[
 e(g(n))=\frac{e(g(n+1))-e(g(n))}{e(g(n+1)-g(n))-1}
\]
and summation by parts bound its sum by $O(\delta^{-1})$: the total
variation of the reciprocal denominator along this arc is
$O(\delta^{-1})$. There are $O_C(M\lambda+1)$ such intervals.
Adding these estimates and inserting $\delta=\sqrt\lambda$ proves the bound.
\end{proof}

\begin{lemma}[spacing, after de Faveri--Pandey]\label{lem:spacingproof}
Let $\phi(x)=\sqrt x+\eta x^{3/2}$. Uniformly for sufficiently small
$|\eta|$, $0<d<1$, and $v\in\mathbb R$,
\[
 \#\{m,n\in(N,2N]:m\ne n,\ |\phi(n/N)-\phi(m/N)-v|\le d\}
 \ll dN^2+N^{2/3}.
\]
\end{lemma}
\begin{proof}
For $d\ge c/N$, monotonicity gives $O(1+dN)$ choices for $n$ at each
$m$, which suffices. Hence suppose $d<c/N$ with $c$ small.
By interchanging $m,n$, take $v\ge0$. If any pair is counted then
$v\gg N^{-1}$ and $v\ll1$. Work on the interval where the inverse below
takes values in a fixed neighborhood of $[1,2]$, and put
$g_v(x)=\phi^{-1}(\phi(x)+v)$. Direct implicit differentiation gives
\[
 g_v'=\frac{\phi'(x)}{\phi'(g_v)},\qquad
 g_v''=\frac{\phi'(x)^2}{\phi'(g_v)}
 \left(\frac{\phi''(x)}{\phi'(x)^2}
       -\frac{\phi''(g_v)}{\phi'(g_v)^2}\right).
\]
Now $\phi'>c_1>0$, $\phi''<-c_1<0$, and
$(\phi''/\phi'^2)'$ is bounded above and below by positive constants
on this fixed interval. The last statement follows at $\eta=0$ from
$\phi''/\phi'^2=-x^{-1/2}$, and persists by continuity.
It follows that $g_v-x\asymp v$, $g_v'-1\asymp v$, and
$-g_v''\asymp v$. Each admissible $m$ must satisfy
\[
 \|f(m)\|\le C dN,\qquad f(u)=N(g_v(u/N)-u/N).
\]
Each $m$ has $O(1)$ corresponding $n$.

If $v\le N^{-1/3}$, the values $f(m)$ run through an interval of length
$O(Nv)$ and are separated by $\gg v$. Thus their number within $CdN$
of an integer is $O((1+Nv)(1+dN/v))=O(Nv+dN^2)$, as $Nv\gg1$.

For $v>N^{-1/3}$ let $b=C'dN+(v/N)^{1/3}$.
If $b$ exceeds a fixed small constant, the trivial bound suffices.
Otherwise periodize a nonnegative smooth bump at scale $b$, equal to at
least one on $[-CdN,CdN]$. Its Fourier coefficients satisfy
$|c_h|\ll_A b(1+b|h|)^{-A}$. Since $|f''|\asymp v/N$,
Lemma~\ref{lem:secondproof} and Fourier expansion give
\[
 \#\{m:\|f(m)\|\le CdN\}
 \ll Nb+\sqrt{Nv}\,b^{-1/2}+\sqrt{N/v}\,b^{1/2}+1.
\]
Here $Nb\ll dN^2+N^{2/3}$,
$\sqrt{Nv/b}\ll N^{2/3}v^{1/3}\ll N^{2/3}$, and
$\sqrt{Nb/v}\ll\sqrt{N/v}\ll N^{2/3}$.
The absolutely convergent smooth Fourier series and the proved
second-derivative bound establish this spacing estimate.
\end{proof}

\section{The curved multiplier and all even energies}\label{app:energyproof}

\begin{proposition}\label{prop:multproof}
Lemma~\ref{lem:mult} holds with an ordinary implied constant in place of
$\prec$. The constants are uniform on fixed support enlargements where
$|\kappa|W$ is sufficiently small.
\end{proposition}
\begin{proof}
Write $\lambda=\xi_m-\xi_n$, $\nu=\chi_m-\chi_n$,
$R_\nu=1+B|\nu|$ and $a_\nu(u)=\beta(u/W)e(\nu r_0(u))$.
On this support,
\[
 |W^j r_0^{(j)}(u)|\ll_j B,\qquad
 r_0''(u)=-2\eta\kappa(1+\kappa u)^{-3},\qquad
 |\nu|\asymp|\lambda|.
\]
The integral part of Lemma~\ref{lem:secondproof}, and the trivial estimate
when $B|\nu|\le1$, show $|\widehat a_\nu(v)|\ll WR_\nu^{-1/2}$.
If $W|v|>C R_\nu$, rescale $u=Wx$ and integrate with
$(2\pi i(\nu W r_0'(Wx)-Wv))^{-1}d/dx$. Its denominator is comparable
to $W|v|$, and every normalized derivative of the nonlinear part is
$O_j(B|\nu|)$. Repeating gives
$|\widehat a_\nu(v)|\ll_J W(1+W|v|)^{-J}$ in this region.

Lemma~\ref{lem:spacingproof} puts $O(N^2/W)$ differences $\lambda$ in
each interval of length $1/W$, since $W\le N^{4/3}$.
On $|u-\lambda|\le CR_\nu/W$, the smallness of $B/W$ gives
\[
 1+B|\lambda|\asymp1+B|u|.
\]
For example $B||\lambda|-|u||\le C(B/W)(1+B|\lambda|)$ and this term
can be absorbed in both directions. The near region has
$O(1+B|u|)$ cells. Multiply their population by
$W(1+B|u|)^{-1/2}$. The result is $O(N^2\sqrt{1+B|u|})$.
On the complement use the separate far estimate just proved, summing
$O(N^2/W)W(1+|j|)^{-J}$ over cells at distance $j/W$ from $u$.
This contributes $O(N^2)$ and proves the statement.
\end{proof}

\begin{proposition}[even energy extension]\label{prop:allenergy}
Under the hypotheses of the curved model, for every fixed integer $j\ge2$
and $1\le D\le W$,
\[
 E_{2j}(S;D)\ll_j
 \left(\frac NZ\right)^{4j-4}
 \left(1+\frac BD\right)^{j-1}E_2(S;D)
 +\frac{ND}{Z^2}r^{2j-1}.
\]
The implicit constants depend on the fixed cutoffs and on $j$.
\end{proposition}
In particular this proves Proposition~\ref{prop:curv}.
\begin{proof}
Set $\mathcal F(u)=\sum_{w\in S}e(wu)$. For each ordered $(2j-1)$-tuple put
$\tau=w_1+\cdots+w_j-w_{j+1}-\cdots-w_{2j-1}$.
Multiply its last-point count by $Z^2$ and use
Lemma~\ref{lem:localproof} on each eligible last point.
One fixed nonnegative $g\in C_c^\infty$, equal to one on $[-2,2]$,
majorizes the resulting local intervals by $g((\tau-u)/D)$.
Since
\[
 D_h'(u)=2\pi i\left[\sum a_n\xi_ne(u\xi_n+r_0(u)\chi_n)
 +r_0'(u)\sum a_n\chi_ne(u\xi_n+r_0(u)\chi_n)\right]
\]
and $|r_0'|\ll B/W\ll1$, we obtain three integrals with fixed
coefficient sequences bounded by an absolute constant.
Before enlargement all integration points lie in $[-1,W+1]$, because
they are local neighborhoods of points of $S$. Insert a fixed
nonnegative cutoff $\beta(u/W)$ equal to one there, supported in
$[-3W,3W]$, and then enlarge the integrals to the whole line.
Each coefficient class and the cutoff remain fixed across tuples and $j$.

For each coefficient class the diagonal after opening the square is
$O_j(NDr^{2j-1})$. For an off-diagonal frequency pair, convolution and
Fourier inversion give the exact identity
\[
 \int g((\tau-u)/D)a_\nu(u)e(\lambda u)\,du
 =D\int\widehat g(Dv)\widehat a_\nu(v-\lambda)e(\tau v)\,dv.
\]
Summing over the $2j-1$ points gives $\mathcal F(v)|\mathcal F(v)|^{2j-2}$.
Proposition~\ref{prop:multproof} therefore bounds the absolute
off-diagonal contribution by
\[
 CN^2\sqrt{1+B/D}\,D\int\Psi(Dv)|\mathcal F(v)|^{2j-1}\,dv,
 \qquad \Psi(v)=(1+v^2)^{-J},
\]
with $J$ fixed sufficiently large. We used
$\sqrt{1+B|v|}\le\sqrt{1+B/D}\sqrt{1+D|v|}$ and rapid decay of
$\widehat g$.

For each fixed $h\ge1$, expansion and Fourier inversion show
\[
 D\int\Psi(Dv)|\mathcal F(v)|^{2h}\,dv\ll_h E_{2h}(S;D).
\]
For the tail bound, $|\widehat\Psi(u)|\ll_A(1+|u|)^{-A}$ follows by integration
by parts. Split the tuple discrepancies into dyadic shells and use
Lemma~\ref{lem:enlarge} on the multiset of $h$-fold sums. The shell at
radius $2^aD$ contributes $O(2^{-(A-1)a}E_{2h}(S;D))$.
H\"older with exponents $2j-2$ and $(2j-2)/(2j-3)$ now yields
\[
 Z^2E_{2j}\ll_j NDr^{2j-1}
 +N^2\sqrt{1+B/D}\,
 E_2^{1/(2j-2)}E_{2j}^{(2j-3)/(2j-2)}.
\]
If the first term is at least half the left side, it gives the stated
diagonal bound. Otherwise divide by the last power of $E_{2j}$ and
raise to $2j-2$. This proves the statement, including all ordered tuples.
\end{proof}

\section{Exact oscillatory integrals and their parameters}\label{app:phaseproof}

\begin{lemma}[two pairs with finitely many derivatives]\label{lem:smoothpairs}
For a fixed compact interval and $f\in C^2$ on a slightly larger interval,
\[
 |f(v_1)+f(v_2)-f(v_3)-f(v_4)|
 \ll\|f\|_{C^2}(|P_v|+|Q_v|).
\]
The statement applies to any fixed parameter derivative of $f$ provided
that derivative has a bounded $C^2$ norm.
\end{lemma}
\begin{proof}
Write a pair as $s/2\pm r$ and put $\rho=r^2$.
For $G(s,\rho)=f(s/2+r)+f(s/2-r)$,
\[
 G_s=\tfrac12(f'(s/2+r)+f'(s/2-r)),\quad
 G_\rho=\tfrac12\int_{-1}^1 f''(s/2+vr)\,dv.
\]
These formulas include $r=0$ by continuity. Extend $f$ smoothly to a
larger fixed interval if necessary and use the segment between the two
points $(s,\rho)$ in the half-plane $\rho\ge0$.
For the two pairs, $s-s'=P_v$ and
$\rho-\rho'=Q_v/2-(s^2-s'^2)/4$.
The mean value theorem proves the estimate.
Applying the same argument to a parameter derivative proves uniformity.
\end{proof}

\begin{lemma}[exact Morse and Fresnel identity]\label{lem:exactmorse}
Let $p$ range in a fixed compact parameter set. Suppose $f(p,x)$ has,
on a fixed interval containing the support of a smooth amplitude $a$,
a unique critical point $x_0(p)$, and $s f_{xx}\ge c>0$ there for one
$s\in\{-1,1\}$. Suppose also that $f$ and $a$ have bounded derivatives
of all fixed orders under consideration. If $Y\ge1$, then exactly
\[
 \int a(p,x)e(Yf(p,x))\,dx
 =Y^{-1/2}e(Yf(p,x_0(p)))\,b(p,Y),
\]
where $(Y\partial_Y)^h\partial_p^\alpha b\ll_{h,\alpha}1$.
The support of the amplitude is allowed to contain $x_0$ or lie away from it.
\end{lemma}
\begin{proof}
The implicit function theorem and differentiation of $f_x(p,x_0)=0$
give bounded derivatives of $x_0$. Taylor's integral formula gives
\[
 v=(x-x_0)\left(s\int_0^1(1-u)f_{xx}(p,x_0+u(x-x_0))\,du\right)^{1/2}.
\]
It has bounded derivatives and an inverse with bounded derivatives;
the identity $2v v_x=s f_x$ and the lower curvature bound show
$v_x\asymp1$, also at $x_0$. The phase becomes
$f(p,x_0)+s v^2$. The transformed amplitude $c(p,v)$ is smooth and
compactly supported in a fixed interval, also after extension by zero.
Fourier inversion and the Gaussian integral give
\[
 b(p,Y)=\frac{e(s/8)}{\sqrt2}
 \int_{\mathbb R}\widehat c(p,\xi)e(-s\xi^2/(4Y))\,d\xi.
\]
One may first insert $e^{-\epsilon v^2}$, evaluate the ordinary Gaussian
integral, and pass to the limit by Schwartz domination on the Fourier
side; this justifies the oscillatory Gaussian identity.
Each $Y\partial_Y$ inserts a polynomial in $\xi^2/Y$.
Parameter derivatives fall on $\widehat c$. Its Schwartz bounds prove
all the stated estimates. The construction uses an interval containing both the support and $x_0$.
It also covers the case $a(x_0)=0$.
\end{proof}

\begin{proposition}[critical value and tuple derivatives]\label{prop:exactphaseproof}
Lemma~\ref{lem:phase} and \eqref{eq:psi} hold. More explicitly, if
$L=NM$, $t=m/M$ is in a fixed interval away from zero, $v_i=w_i/L$
are bounded and $\gamma=\kappa L$, then
\[
 \Phi(t,x)=\frac Q{4Lt}-\frac{Px}{2}+C(t,x),\qquad
 |\partial_t^a\partial_x^b C|\ll_{a,b}|\eta|(|P|+F),
\]
and
\[
 C_{xx}=\frac{3\eta Q}{2Lt}-3\eta Px+E(t,x),\qquad
 |\partial_t^a\partial_x^b E|
 \ll_{a,b}(|\gamma|+|\eta|)|\eta|(|P|+F).
\]
All estimates are on fixed compact $x$ intervals, with their smallness
constants chosen after those intervals.
\end{proposition}
\begin{proof}
The pair phase before Poisson integration is
\[
 y\{\Delta u^{1/2}+\eta_{\rm eff}\Delta u^{3/2}-u\},\quad
 y=Nm,\quad \Delta=(z-w)/y.
\]
Writing $r=\sqrt u$ in its critical equation gives
$3\eta_{\rm eff}\Delta r^2-2r+\Delta=0$.
The root remaining in a fixed positive interval as $\eta\to0$ is
$r=\Delta/(1+\sqrt{1-3\eta_{\rm eff}\Delta^2})$.
The other root is outside the fixed integration support. On the
stationary range $\Delta$ is positive and bounded above and below,
and $|\eta_{\rm eff}\Delta^2|<1/4$ follows from smallness of $|\eta|$.
Substitution, using $\frac34\tau\Omega^2-\Omega+1=0$, gives exactly
\[
 \Theta=\frac{y\Delta^2}{12}\Omega(4-\Omega).
\]
The analytic expansion at $\tau=0$ is
\[
 \Theta=\frac{y\Delta^2}{4}
       +\frac{y\eta_{\rm eff}\Delta^4}{8}
       +y\eta_{\rm eff}^2\Delta^6\Psi(\tau),
 \qquad \Psi(0)=\frac9{64}.
\]
This identity keeps the full analytic remainder function in the phase.

At $z=Ltx$ its leading term sums to $Q/(4Lt)-Px/2$.
The next term sums to
\[
 \frac{\eta L}{8t^3(1+\gamma tx)}
 \sum_i\epsilon_i\frac{(tx-v_i)^4}{1+\gamma v_i}.
\]
The remainder is $\eta^2L$ times the signed two-pair difference of a
uniformly smooth function of $v_i$, also after any fixed $t,x$
derivatives. Lemma~\ref{lem:smoothpairs} bounds it by
$O_{a,b}(\eta^2(|P|+F))$. The same lemma bounds the quartic term by
$O_{a,b}(|\eta|(|P|+F))$.
Replacing $\gamma$ by zero in its second $x$ derivative costs
$O_{a,b}(|\gamma\eta|(|P|+F))$ by the mean value theorem in $\gamma$
and the two-pair estimate. At $\gamma=0$ direct differentiation gives
$3\eta Q/(2Lt)-3\eta Px$. This proves every claimed derivative bound.
\end{proof}

\begin{proposition}[tuple Morse formula]\label{prop:tuplemorseproof}
Lemma~\ref{lem:morse} holds. The nonstationary and small-$H$ statements
used in Proposition~\ref{prop:kernel} also hold with bounded derivatives
of every fixed required order.
\end{proposition}
\begin{proof}
Let $H=|\eta|F$. When $|P|>C(1+H)$, the first $x$ derivative has
magnitude $\asymp|P|$. Factor $e(Q/(4Lt))$ and integrate by parts in
$x$. Every fixed $t$ derivative costs at most a fixed power of
$|P|+H$; choosing that many extra integrations gives
$O_{a,J}((1+|P|)^{-J})$ for the remaining amplitude and its derivatives.
For $H\le1$ and $|P|\le C(1+H)$ all derivatives of $\psi$ are bounded,
so $e(\psi)$ is a bounded smooth amplitude.

Suppose $H\ge1$ and $|P|\le CH$. The preceding proposition gives
\[
 \psi_{xx}/H=\operatorname{sgn}(\eta Q)\frac3{2t}
             +O(|\eta|+|\kappa|L)
\]
for positive $t$; the sign is adjusted for negative $t$.
Moreover $\psi_x/H$ is an affine function of $x$ with bounded intercept
and slope separated from zero, plus $O(|\eta|+|\kappa|L)$ in every
fixed derivative on a fixed interval. For the intercept one uses
$|\sum\epsilon_iw_i^3|\ll L|Q|+L^2|P|$, from
Lemma~\ref{lem:smoothpairs}. Choose a fixed larger interval containing
the amplitude support and all roots of these affine functions, then
choose the small constants. The derivative has opposite signs at the
ends, so there is exactly one critical point and its derivatives are
bounded. Apply Lemma~\ref{lem:exactmorse} with $Y=H$ and $f=\psi/H$.
Here $H$ is fixed as $t$ varies. It yields
$H^{-1/2}b(t)e(Q/(4Lt)+\psi(t,x_0(t)))$, with
$b^{(a)}\ll_a1$ and $(d/dt)^a\psi(t,x_0(t))\ll_aH$.
The construction includes critical points outside the amplitude support.
\end{proof}

\section{A complete signed duality calculation}\label{app:dualityproof}

Put $L=NM$. We write the bound as a sum over height blocks. Partition $S$ into
intervals of length $CL$, independently of $m$, where $C$ is fixed
sufficiently large. For each outer block $S_j$ and neighboring inner
block $S_{j+h}$, $h\in\{-1,0,1\}$, translate by its fixed origin $u_j$.
Let $\mathcal B_{M,j,h}$ be the translated inner set. Then
$|w|\ll L$ in $\mathcal B_{M,j,h}$. The parameters become
\[
 \eta_j=\eta(1+\kappa u_j)^{-2},\qquad
 \kappa_j=\kappa(1+\kappa u_j)^{-1}.
\]
The common phase at $u_j$ cancels in the pair difference, since
$h(u_j+v)-h(u_j)=\eta_jv/(1+\kappa_jv)$. Thus the translation can be
performed after the global duality step with its coefficients preserved.

\begin{proposition}[signed blocks]\label{prop:dualityproof}
There are smooth amplitudes in the two fixed Sobolev classes and in
each sign of $m$ for which
\begin{equation}\label{eq:preciseduality}
 r^2Z^2\ll N^2r+
 \sum_{M,h,\pm,\ell}N^{7/4}M^{1/2}
 \sum_j |S_j|^{3/4}\mathcal T_{M,j,h,\pm,\ell}^{1/4},
\end{equation}
where $\ell=0,1$, $\mathcal T\ge0$, and
\[
 \mathcal T_{M,j,h,\pm,\ell}
 =\sum_{\mathbf w\in\mathcal B_{M,j,h}^4}
 b_{w_1}b_{w_2}\overline{b_{w_3}b_{w_4}}
 \sum_{m\text{ in the fixed sign}}\int A_{\ell,\mathbf w}(m/M,x)
 e(\Phi(m/M,x))\,dx.
\]
All normalized derivatives of the amplitudes are bounded. In particular
$\mathcal T\le\sum_{\mathbf w}|K_{M,j,h,\pm,\ell}(\mathbf w)|$.
The number of neighbors and signs is bounded independently of $N,W,r$.
\end{proposition}
\begin{proof}
Choose $b_w$ of modulus one so that $b_wD_h(w)=|D_h(w)|$.
Cauchy's inequality only in the $n$ variable gives
\[
 r^2Z^2\le N\sum_n\varphi(n/N)
 \left|\sum_{w\in S}b_w
 e(w\sqrt{n/N}+h(w)(n/N)^{3/2})\right|^2,
\]
where $\varphi\ge1$ on $[1,2]$ and has fixed compact positive support.
Poisson summation changes the right side exactly to
\[
 N^2\sum_{z,w\in S}b_z\overline b_w\sum_{m\in\mathbb Z}
 \int\varphi(u)e((z-w)\sqrt u+(h(z)-h(w))u^{3/2}-Nm u)\,du.
\]
Poisson follows, for this smooth compact support, by periodizing the
summand and taking its absolutely convergent Fourier series.
For $m=0$ the derivative is comparable to $z-w$ unless $z=w$;
integration by parts and $1$-spacing give $O(N^2r)$.

For $m\ne0$, stationary support requires $\Delta=(z-w)/(Nm)$ in a
fixed compact subinterval of $(0,\infty)$. Choose a smooth cutoff in
$\Delta$ equal to one on a neighborhood of all such ratios. Off that
cutoff, repeated integration by parts bounds the integral by
$O_A((1+|z-w|+N|m|)^{-A})$.
For each fixed $z$, $1$-spacing implies
$\sum_w(1+|z-w|+N|m|)^{-A}\ll(1+N|m|)^{1-A}$.
Summing in $m$ bounds this part by $O(N^2r)$.

On the cutoff use the exact identity
$h(z)-h(w)=\eta(z-w)/((1+\kappa z)(1+\kappa w))$.
The pair phase divided by $|Nm|$ has one nondegenerate critical point
in a fixed positive enlargement of the $u$ support. Its second
derivative has constant sign there, bounded away from zero. Apply
Lemma~\ref{lem:exactmorse} with $Y=|Nm|$.
The integral is \emph{exactly}
\[
 (N|m|)^{-1/2}a_m(z,w)e(\Theta_m(z,w)).
\]
This exact representation keeps the full stationary integral. Smooth cutoffs in
$\Delta$ are included in $a_m$. Its normalized derivatives in
$m/M,z/L,w/L$ are bounded: derivatives of $Y$ are controlled by the
$Y\partial_Y$ estimates, and derivatives of the remaining parameters
cost bounded powers of $|\kappa|L$.

Fix one dyadic scale, neighbor and sign. The contribution is bounded by
\[
 N^{3/2}M^{-1/2}
 \sum_j\sum_m\omega(m/M)\sum_{z\in S_j}|f_{m,0}(z-u_j)|,
\]
where $\omega$ is a fixed nonnegative smooth dyadic weight and
\[
 f_{m,0}(z)=\sum_{w\in\mathcal B_{M,j,h}}b_w a_{m,0}(z,w)e(\Theta_m(z,w)).
\]
Weighted H\"older in $(m,z)$ gives the factor
$(M|S_j|)^{3/4}$ and the fourth root of
\[
\sum_m\omega(m/M)\sum_{z\in S_j-u_j}|f_{m,0}(z)|^4.
\]
By Lemma~\ref{lem:localproof} the latter is bounded by
\[
 \sum_m\omega(m/M)\int\chi_j(z/L)
 (|f_{m,0}(z)|^4+|f_{m,1}(z)|^4)\,dz,
 \quad f_{m,1}=\partial_zf_{m,0}.
\]
The derivative amplitude is
$a_{m,1}=\partial_z a_{m,0}+2\pi i(\partial_z\Theta_m)a_{m,0}$.
At the critical point the envelope identity gives
$\partial_z\Theta_m=\sqrt{u_0}+h'(z)u_0^{3/2}=O(1)$,
with all normalized derivatives bounded. Thus both classes qualify.

Set $z=Nmx$. The Jacobian is $N|m|=L|m|/M$.
Its smooth factor $|m|/M$, and the cutoff $\chi_j((m/M)x)$, belong in
$A_{\ell,\mathbf w}$. Opening each fourth power before taking absolute
values produces exactly $\mathcal T$ above. It is nonnegative because
it was the integral of a fourth power against nonnegative weights.
The remaining external factor is
\[
 N^{3/2}M^{-1/2}(M|S_j|)^{3/4}L^{1/4}
 =N^{7/4}M^{1/2}|S_j|^{3/4}.
\]
Both $b_w$ and the translation remain fixed across $m$. This proves
\eqref{eq:preciseduality} and accounts for both signs, both derivative
classes, all neighboring blocks and all nonstationary modes.
\end{proof}

\begin{corollary}[the energy bound]\label{cor:blockconsumer}
Define
\[
 \mathfrak K_M=\sum_j\sum_{\mathbf w\in\mathcal B_{M,j,h}^4}
 |K_{M,j,h,\pm,\ell}(\mathbf w)|,
\]
separately for each finite class, and
$\cJ=\max_{M,h,\pm,\ell}(M/\Mst)^3M^{-1}\mathfrak K_M$.
Then \eqref{eq:rJ} and \eqref{eq:Jbound} hold with this definition.
\end{corollary}
\begin{proof}
H\"older in $j$ bounds the block sum in
\eqref{eq:preciseduality} by $r^{3/4}\mathfrak K_M^{1/4}$.
Dyadic summation loses only $O(\log N)$.
For each fixed $h$, the inner blocks partition $S$; hence
\[
 \sum_jE_4(S_{j+h};D)\le E_4(S;D),\qquad
 \sum_j|S_{j+h}|E_2(S_{j+h};H_0)\le rE_2(S;H_0).
\]
The discrepancies $P$ are translation invariant. The local $Q$, hence
$F$, is calculated after its block's fixed translation; the small-$F$
proof is applied there before going back to the translation invariant
close-pair counts. Each block keeps its own translated $Q$.
The pointwise kernel bounds are uniform over the translated parameters,
so energy bounds can be applied after summing the block tuple counts.
The remainder of the calculation is exactly the algebra in
Proposition~\ref{prop:model}, with $\mathfrak K_M$ in place of a single
block kernel. This proves both stated inequalities.
\end{proof}

\section{The reduction to the zeta function}\label{app:transferproof}

\begin{lemma}[partial summation and interval completion]\label{lem:completionproof}
The first phase truncation in Lemma~\ref{lem:atkinson} changes a partial
sum by at most a constant times the maximum of the partial sums with
the truncated phase. For finitely many blocks and any phase $p_j(w,n)$,
the sum over sample points of these maxima is bounded by
$C\log(2N)$ times the sum of full sums after multiplying every
coefficient by one common factor $e(hn/(2N))$ for a suitable integer $h$.
\end{lemma}
\begin{proof}
For $A(x)=\sum_{N<n\le x}c_n$ and $v(n)=e(r_t(n))$, Abel summation gives
\[
 \sum_{N<n\le y}c_nv(n)=A(y)v(y)-\int_N^y A(x)v'(x)\,dx.
\]
The factor is at most $1+2\pi\int_N^{2N}|r_t'(x)|dx=O(1)$ when
$N^{5/2}T^{-3/2}\ll1$.
For interval completion, expand the indicator of $1\le n-N\le\Delta$
on $\mathbb Z/(2N)\mathbb Z$. Its coefficient at $h$, $-N<h\le N$,
is bounded by $C/(1+|h|)$, directly from the finite geometric sum.
Thus each maximum is bounded by
\[
 C\sum_{-N<h\le N}\frac1{1+|h|}
 \left|\sum_{n\sim N}a_{j,n}e(hn/(2N))e(p_j(w,n))\right|.
\]
Sum this inequality over the blocks and sample points. Choosing the
$h$ with the largest resulting sum proves the statement.
This argument applies to every phase $p_j(w,n)$. Nonintegral endpoints are
handled by their integer parts or by a cyclic group of length $\asymp N$.
\end{proof}

\begin{proposition}[refined pair count from sparse fourth moments]
\label{prop:pairproof}
Assume the sparse fourth-moment input \eqref{eq:basefourth} and the
global fourth moment. Let $t_1,\ldots,t_R$ be $2G$-separated, and assume
\[
\int_{t_j-G}^{t_j+G}|\zeta(1/2+iu)|^4du\ge A\ge1.
\]
Then for all $H\ge1$,
\[
 E_2(\{t_j\};H)\prec R+RH/A+T\sqrt R/A^{3/2}.
\]
\end{proposition}
\begin{proof}
Disjointness and the fourth moment give $RA\prec T$.
For $H<G$ there are only diagonal pairs; for $H\ge T$ use
$R^2\prec RT/A$. Otherwise partition into cells of length $H$ with
occupancies $m_j$. Adjacent-cell Cauchy gives $E_2\le3\sum_jm_j^2$.
It suffices to treat cells of occupancy in $[m,2m)$, their number being
$J$, since there are $O(\log T)$ scales. We have $Jm\le R$.
Fix a small $\delta>0$.
If $mA\le T^\delta\max(H,\sqrt T)$, then
\[
 Jm^2\le Rm\ll T^\delta(RH/A+R\sqrt T/A)
 \prec T^\delta(RH/A+T\sqrt R/A^{3/2}).
\]
In the other case set $L=\min(mAT^{-\delta},T)>\max(H,\sqrt T)$.
The $G$-neighborhoods of the points in these $J$ cells can be covered
by $O(J)$ intervals of length $L$, with a bounded number of disjoint
families. Endpoint pieces can be handled on enlarged dyadic height
intervals. Therefore the sparse fourth moment gives
\[
 JmA\prec JL+T\sqrt{J/L}.
\]
Choose its arbitrary-small-power loss smaller than $\delta/2$ to
absorb $JL$. Squaring and using $mA\prec T$ gives
$Jm^3\prec T^{2+O(\delta)}A^{-3}$.
Thus $Jm^2=\sqrt{Jm}\sqrt{Jm^3}\prec
T^{1+O(\delta)}\sqrt R A^{-3/2}$.
As $\delta$ is arbitrary the proposition follows.
\end{proof}

\begin{proposition}[clustering from the stated zeta theorems]
\label{prop:clusterproof}
Proposition~\ref{in:dfp}, with $1/10<\sigma\le1/8$, follows from
\eqref{eq:baselocal}, \eqref{eq:baseatkinson},
\eqref{eq:basefourth} and the global fourth moment. The same conclusion
holds for $1/10<\sigma<1/6$ when the corresponding parameter conditions
for the curved model are satisfied.
\end{proposition}
\begin{proof}
Start with any $1$-separated zeta large-value set of size $R$, and put
$G=V^2(\log T)^{-100}$. Partition into cells of length $2G/3$.
Dyadic occupancy selection and one residue class modulo four give
a $2G$-separated set $\mathcal S$ of representatives, of size
$\succ R/F_c$, each representing $\asymp F_c\ge1$ original points.
By the local inequality \eqref{eq:baselocal}, their $G$-neighborhoods
carry second and fourth moments at least $\succ F_cV^2$ and
$\succ F_cV^4$, respectively.

The term $G\log T$ in \eqref{eq:baseatkinson} is smaller than the
second-moment lower bound. Select one dyadic $N$ for at least a
logarithmic fraction of the representatives. Then
\[
 T^{1/3}\ll N\prec T/V^4,\qquad
 \max_\Delta|S_t(\Delta,N)|\succ F_c(TN)^{1/4}.
\]
Since $\sigma>1/10$, the strict margin gives $N<T^{3/5}$ after all
logarithmic factors. Divide the ordinates into blocks of length
$\asymp W\sqrt{T/N}$. There are $O(1+\sqrt{TN}/W)$ blocks.
The coordinate in Lemma~\ref{lem:atkinson} sends each to $[0,CW]$;
fixed constants in $W$ may be chosen to obtain $[0,W]$.
The exact rational height increment, partial summation, the divisor
bound and Lemma~\ref{lem:completionproof} give a sum of absolute full
model values at least $\succ R(TN)^{1/4}$.

Discard values below a sufficiently small arbitrary-small-power
multiple of $F_c(TN)^{1/4}$; their total is less than half this last
sum. Dyadic selection among the remaining values gives $Z$ and sets
of total size $R_1$ with $Z\succ F_c(TN)^{1/4}$ and
$ZR_1\succ R(TN)^{1/4}$. This argument uses the summed inequality after selecting one common
Fourier index; the set and its level $Z$ are then selected together.

The transformed spacing is $\gg G\sqrt{N/T}$. Color the points of
each block in their order into $O(1+G^{-1}\sqrt{T/N})$ classes so each
class is $1$-separated. Because $N\prec T/V^4$ and $W\le\sqrt{TN}$,
the total number of classes is $\prec T/(V^2W)$.
Each point occurs once. For two points in a single class,
$|w-w'|\le D$ implies $|t-t'|\ll D\sqrt{T/N}$.
Apply Proposition~\ref{prop:pairproof} to the entire selected set of
ordinates $t$ with $A\succ F_cV^4$.
This proves \eqref{eq:paircount} for every $D\ge1$, including its
large-radius regime. The $1$-spacing bound is independent and remains
available simultaneously.
\end{proof}

\section{The upper-amplitude range}
\label{app:upperproof}

\begin{proposition}\label{prop:curvatureline}
Subject to the zeta inputs of Proposition~\ref{prop:clusterproof},
\[
 \cN(T,T^\sigma)\prec T^{(94-500\sigma)/67},
 \qquad 1/8\le\sigma\le5/38.
\]
\end{proposition}
This statement follows from the direct oscillatory-integral estimates.
\begin{proof}
Proposition~\ref{prop:tuplemorseproof}, applied at each fixed $m$, gives
$|I_m|\ll D^{-1/2}(1+D/U)^{-A}$ for every $M$ and every tuple.
Use this pointwise bound on all tuples, including the small-$F$ range. The signed-block calculation and the curved energy bound give
\[
 r\prec N^2/Z^2+N^{8/5}W^{3/5}Z^{-12/5}A_0(S)^{1/5}
 +N^{5/2}W^{3/2}Z^{-5}U^{1/4}
 +B^{1/4}N^2W^{3/4}Z^{-3}.
\]
The weighted pair argument with weight $D^{-1/2}$ is valid also for
$k=0$. Thus the reduction to the zeta function gives exactly the terms
$B_1,\ldots,B_6$ with
$\lambda=\mu=1$ and $k=0$. Put
\[
 \omega(\sigma)=\frac{36+48\sigma}{67},\qquad
 q=\frac{\omega(\sigma)}{1-4\sigma},\qquad W=N^q.
\]
Here $84/67\le q\le4/3$. All six exponents increase with
$\nu=\log_TN$. At $\nu=1-4\sigma$, terms 4 and 5 equal
$E_c=(94-500\sigma)/67$. The margins for terms 1, 2, 3 and 6 are
\[
 \frac{4(21\sigma-1)}{67},\quad
 \frac{5-38\sigma}{134},\quad
 \frac{9-55\sigma}{67},\quad
 \frac{5-38\sigma}{268},
\]
which are nonnegative throughout the closed interval.
For the model ranges and pair absorption the exponents of $T$ in
$|\kappa|W$, $|\eta|N^2/W$, $|\eta|N^{1/8}$,
$N^{5/2}T^{-3/2}$ and $W\sqrt{N/T}V^{-4}$ are, respectively,
\[
 \frac{-31+182\sigma}{67},\quad
 \frac{98-852\sigma}{67},\quad
 \frac18-\frac92\sigma,\quad
 1-10\sigma,\quad
 \frac{36-354\sigma}{67}.
\]
They are strictly negative at both endpoints. These quantities also
have positive $\nu$ slopes for fixed $q$, except the separately treated
$\sqrt{T/N}$ term, which is at most $T^{1/3}<V^4$.
The endpoint $q=4/3$ is permitted in the spacing estimate.
This proves the proposition. Its contribution to $M_6$ is at most
$T^{327/268+\varepsilon}$, attained at $\sigma=1/8$.
\end{proof}

\section{Local zeta values and consequences of sparse fourth moments}\label{app:localzeta}

\begin{proposition}[local values from harmonic measure]\label{prop:localzetaproof}
For every fixed $p>0$,
\begin{equation}\label{eq:baselocal}
 |\zeta(1/2+it)|^p\ll_p
 1+\log(2+|t|)\int_{t-1/3}^{t+1/3}|\zeta(1/2+iu)|^p\,du.
\end{equation}
\end{proposition}
\begin{proof}
We give the localization argument, including the analytic estimates it
requires. The elementary functional equation gives, uniformly for
$0\le u\le1/\log T$ and $\tau\asymp T$,
\begin{equation}\label{eq:localreflection}
 |\zeta(1/2-u+i\tau)|\ll |\zeta(1/2+u+i\tau)|.
\end{equation}
Here is one derivation of the facts needed in this statement. Periodizing
the Gaussian, whose Fourier transform is $x^{-1/2}e^{-\pi\xi^2/x}$,
gives $\theta(x)=x^{-1/2}\theta(1/x)$. Mellin integration, initially
for $\Re s>1$, and splitting at $x=1$, then gives
\[
 \pi^{-s/2}\Gamma(s/2)\zeta(s)
 =-\frac1s-\frac1{1-s}
 +\frac12\int_1^\infty(\theta(x)-1)
       \bigl(x^{s/2}+x^{(1-s)/2}\bigr)\,\frac{dx}{x}.
\]
The right side continues meromorphically and is unchanged by
$s\mapsto1-s$. The gamma quotient therefore has modulus
$\asymp(1+|\tau|)^u$ at $s=1/2-u+i\tau$, which proves
\eqref{eq:localreflection}. The gamma estimate follows
by applying Euler summation twice to the logarithmic gamma product:
on $\Re z>0$ it gives
\[
 \log\Gamma(z)=(z-\tfrac12)\log z-z+c+\frac1{12z}
 -\frac12\int_0^\infty\frac{B_2(\{x\})}{(x+z)^2}\,dx.
\]
The integral is $O(|\Im z|^{-1})$ when $\Re z$ is in a fixed positive
compact interval; take real parts in the two gamma factors.
Conjugation accounts for the sign of $\tau$. Also Euler summation in
the defining Dirichlet series, with a cutoff $X\asymp1+|\tau|$, gives
\[
 \zeta(s)=\sum_{n\le X}n^{-s}+\frac{X^{1-s}}{s-1}
          +O(X^{-\Re s}+|s|X^{-\Re s}/\Re s)
\]
for $\Re s>0$ away from $s=1$. It gives the polynomial growth used
below directly from Euler summation.

We prove the bound for $T\le t\le2T$ with $T\ge10$. On bounded intervals, continuity gives the result.
Set $s_0=1/2+it$, $\delta=1/\log T$, and choose a fixed sufficiently
large $C=C_p$. On the strip $1/2\le\Re s\le3/2$ put
\[
 G(s)=\frac{s-1}{s+1}\zeta(s)
 \exp\bigl(C\log T\bigl[(s-s_0)^2-(s-s_0)\bigr]\bigr).
\]
The pole is canceled. On either boundary the exponential has modulus
$\exp(-C\log T\,|\Im s-t|^2)$. On the right boundary $G$ is bounded
by an absolute constant. In the small rectangle
$0\le\Re(s-s_0)\le\delta$, $|\Im(s-s_0)|\le\delta$,
the other factors show $|\zeta(s)|\ll_C|G(s)|$.

The harmonic-measure kernel for the left boundary of the strip of width
one is
\[
 K_u(v)=\frac{\sin(\pi u)}{2(\cosh(\pi v)-\cos(\pi u))},
 \qquad 0<u<1.
\]
This formula follows by mapping $0<\Re z<1$ to the upper half-plane
by $e^{\pi i z}$ and changing variables in the Poisson kernel.
It is positive, has integral $1-u\le1$, and is bounded by
$C u/(u^2+v^2)$ for $u\le1/2$. Subharmonicity of $|G|^p$ and the
maximum principle on finite rectangles, followed by a limit, give
\[
 |\zeta(1/2+u+i\tau)|^p
 \ll_{p,C}1+\int_{\mathbb R}K_u(\tau-v)
 |\zeta(1/2+iv)|^p T^{-pC(v-t)^2}\,dv
\]
in the small rectangle. The horizontal boundary terms vanish because
of Gaussian decay and the polynomial bound just obtained.

Apply the disk submean inequality to $|\zeta|^p$ at $s_0$ with radius
$\delta$. Reflect the left half using \eqref{eq:localreflection},
and enlarge the right half to the small rectangle. For every $v$,
\[
 \int_0^\delta\int_{t-\delta}^{t+\delta}K_u(\tau-v)\,d\tau\,du
 \le\delta.
\]
After division by the disk area, the bound is
\[
 |\zeta(s_0)|^p\ll_p
 1+\delta^{-1}\int_{\mathbb R}
 |\zeta(1/2+iv)|^pT^{-pC(v-t)^2}\,dv.
\]
For $|v-t|>1/3$, the polynomial bound and sufficiently large $C_p$
make this integral $O(\delta)$, including $|v|$ far larger than $T$.
For the remaining interval discard the Gaussian. This proves
\eqref{eq:baselocal}. The disk submean inequality itself follows from
the circle mean inequality for the subharmonic function $|\zeta|^p$
by integrating the radius, including disks containing zeros.
\end{proof}

\begin{proposition}[fourth and twelfth moments from one sparse input]\label{prop:momentbaseproof}
Assuming only the sparse-fourth-moment estimate
\eqref{eq:basefourth} among the deep zeta moment inputs, one has
\[
 M_4(T)\prec T,\qquad M_{12}(T)\prec T^2,\qquad
 \cN(T,V)\prec\min(TV^{-4},T^2V^{-12})\quad(V\ge1).
\]
\end{proposition}
This derivation obtains the classical twelfth moment from the
sparse fourth moment. Attribution remains with the original authors;
this implication is Iwaniec's route~\cite{Iwa80}.
\begin{proof}
Take $J=1,L=T$ in \eqref{eq:basefourth} and sum dyadic height intervals
to obtain $M_4(T)\prec T$. Proposition~\ref{prop:localzetaproof}
with $p=4$ and disjoint neighborhoods of a $1$-separated set gives
$RV^4\prec T$; the added constant is absorbed for $V$ sufficiently
large, and the bounded-$V$ range is trivial.

Fix $\rho>0$, with all occurrences of $\prec$ in this proof taken at
a loss much smaller than $\rho$. If $V^4\le T^{1/2+\rho}$, the fourth
moment gives
$R\prec TV^{-4}\le T^{2+2\rho}V^{-12}$.
Otherwise put $L=cV^4T^{-\rho}$, with a fixed positive constant $c$.
If necessary increase the lower threshold by a constant so that
$L\ge\sqrt T$; the transition is covered by the preceding case.
When $R\ge1$, the just-proved fourth-moment estimate and the local
inequality give $V^4\ll T^{1+\rho/2}$, hence $L<T$ for large $T$.

Cover the disjoint neighborhoods $[t_j-1/3,t_j+1/3]$ by grid intervals
of length $L$. At most $2R$ grid intervals are needed. Boundary
intervals at $T,2T$ may be enlarged or shifted; a fixed number of
families in comparable dyadic height ranges handles them. Thus
\eqref{eq:basefourth} and the local lower mass yield
\[
 RV^4\prec RL+T\sqrt{R/L}.
\]
The $RL$ term is absorbed, since $L\ll V^4T^{-\rho}$.
Squaring the remaining inequality gives
$R\prec T^2/(LV^8)\ll T^{2+\rho}V^{-12}$.
As $\rho$ is arbitrary this proves the discrete twelfth-moment bound.

Finally the measure of each amplitude layer is at most a constant times
the maximal $1$-separated cardinality: choose a maximal separated
subset, whose unit neighborhoods cover that layer. Integrate in
dyadic $V$ and use the polynomial pointwise bound from the preceding
proof to ensure only $O(\log T)$ layers. This gives $M_{12}(T)\prec T^2$.
\end{proof}

\section{The derivative estimate from a precise mean-value input}\label{app:hbderivative}

We rederive the Heath-Brown estimate used here, keeping its attribution
to \cite{HB17}. Its sole nonelementary input in this appendix is
\begin{equation}\label{eq:vmvtinput}
 J_{s,q}(H):=\int_{[0,1]^q}\left|\sum_{h=1}^{\lfloor H\rfloor}
 e(\alpha_1h+\cdots+\alpha_qh^q)\right|^{2s}d\boldsymbol\alpha
 \ll_{q,\epsilon}H^{s+\epsilon},\qquad s=q(q+1)/2.
\end{equation}
This is the optimal Vinogradov mean-value theorem \cite{BDG16,Woo16}.

\begin{lemma}[coefficient clustering]\label{lem:hbcluster}
Suppose $d\ge3$ and $\lambda\le f^{(d)}\le A\lambda$ on an interval
of length $M$, with $A$ fixed. Put $H=\lfloor(A\lambda)^{-1/d}\rfloor$.
If $A\lambda$ is a sufficiently small constant depending on $d$, then
\[
 \mathcal P:=\#\left\{m,n:\left\|\frac{f^{(j)}(m)-f^{(j)}(n)}{j!}\right\|
 \le2H^{-j}\ (1\le j<d)\right\}
 \ll_{A,d}(M+\lambda M^2+\lambda^{-2/d})\log(2M).
\]
\end{lemma}
\begin{proof}
Keep the last two derivative conditions to obtain an upper bound. Let
$g_1=f^{(d-2)}/(d-2)!$, $g_2=f^{(d-1)}/(d-1)!$, and periodize
triangular functions of widths $B=4H^{2-d}$ and $C=4H^{1-d}$.
Their product $P(x,y)$ is at least $1/4$ on the counted differences,
vanishes unless $\|x\|\le B,\|y\|\le C$, and has nonnegative Fourier
coefficients $BC\,\operatorname{sinc}(rB)^2\operatorname{sinc}(sC)^2$.
Smallness of $\lambda$ ensures $B,C<1/2$.

Partition the integer interval into $K\asymp1+A\lambda M$ consecutive
pieces whose lengths make $|g_2(m)-g_2(n)|<1/2$ within each piece.
Fourier expansion and $|\sum_{i\le K}z_i|^2\le K\sum_i|z_i|^2$ show
\[
 \mathcal P\ll K\sum_i\sum_{m,n\text{ in piece }i}
 P(g_1(m)-g_1(n),g_2(m)-g_2(n)).
\]
Within a piece the $g_2$ condition is an ordinary real difference,
and monotonicity gives $|m-n|\ll\lambda^{-1}H^{1-d}\asymp H$.
The diagonal costs $M$. For $m=n+h$, $1\le h\ll H$, the derivative
of $g_1(x+h)-g_1(x)$ lies between fixed positive multiples of
$\lambda h$. A monotone function with derivative in $[\mu,A'\mu]$
takes integer arguments within distance $b$ of an integer at most
$O((1+\mu M)(1+b/\mu))$ times: count the integers in its range and
the length of each inverse interval. Hence the off-diagonal
count, even after dropping the piece restriction, is at most
\[
 \sum_{h\ll H}(1+\lambda Mh)(1+B/(\lambda h))
 \ll H+\lambda MH^2+(B/\lambda)\log(2H)+MBH
 \ll M+H^2\log(2M).
\]
Here $\lambda\asymp H^{-d}$ and $d\ge3$; if $H>M$ the trivial
$M^2\le H^2$ bound suffices. Multiply by $K$ and use
$\lambda M H^2\ll M$. This proves the bound.
\end{proof}

\begin{proposition}[Heath-Brown from \eqref{eq:vmvtinput}]
\label{prop:hbderivativeproof}
The mean-value bound \eqref{eq:vmvtinput} with $q=d-1$ implies
\eqref{eq:basederivative} for $d\ge3$.
\end{proposition}
\begin{proof}
Write $s=d(d-1)/2$ and first consider an integer interval of length $M$.
Conjugation makes $f^{(d)}$ positive. If $\lambda$ is bounded below,
the claimed first term bounds the sum trivially. If $H\ge M$, the
third claimed term does so. Otherwise $1\ll H<M$ and averaging
the $H$ translates of the sum, with end intervals costing $O(H)$, gives
\[
 \left|\sum_ne(f(n))\right|
 \ll H+H^{-1}\sum_n\max_{x\le H}
 \left|\sum_{h\le x}e\left(\sum_{j=1}^{d-1}\frac{f^{(j)}(n)}{j!}h^j\right)\right|.
\]
Taylor's remainder was removed by Abel summation: its derivative is
$O_d(A\lambda H^{d-1})=O_d(H^{-1})$.
To avoid a point-dependent maximum, use Fourier completion on
$\mathbb Z/(2H)\mathbb Z$, as in Lemma~\ref{lem:completionproof}.
This costs $O(\log(2H))$ and adds a single common linear coefficient
to every polynomial, preserving all coefficient differences.

Let $\nu(\boldsymbol\alpha)$ count the $n$ for which all the polynomial
coefficients are within $H^{-j}$ of $\alpha_j$ modulo one. Perturbing
coefficients by these amounts has total phase variation $O_d(1)$
on $[0,H]$. Abel summation and averaging over the coefficient box
therefore bound the preceding $n$ sum by
\[
 C_dH^s\left(I(H)+H^{-1}\int_0^H I(x)\,dx\right),\qquad
 I(x)=\int_{[0,1]^{d-1}}\nu(\boldsymbol\alpha)
 \left|\sum_{h\le x}e(\alpha_1h+\cdots+\alpha_{d-1}h^{d-1})\right|
 d\boldsymbol\alpha.
\]
Box volumes and overlaps give
$\int\nu\ll MH^{-s}$ and $\int\nu^2\ll\mathcal P H^{-s}$.
H\"older with three factors, of reciprocal exponents
$1-1/s$, $1/(2s)$ and $1/(2s)$, gives
\[
 I(x)\le(\int\nu)^{1-1/s}(\int\nu^2)^{1/(2s)}
 J_{s,d-1}(H)^{1/(2s)}.
\]
The monotonicity in $H$ follows directly by expanding the mean value
as a number of solutions. Applying \eqref{eq:vmvtinput} and collecting
the powers of $H$, which cancel, yields
\[
 \left|\sum_ne(f(n))\right|
 \ll_{d,A,\epsilon}H+M^{1-1/s+\epsilon}\mathcal P^{1/(2s)}.
\]
Insert Lemma~\ref{lem:hbcluster}. Since $2s=d(d-1)$, this is
\[
 \ll M^\epsilon\left(\lambda^{-1/d}
 +M^{1-1/[d(d-1)]}+M\lambda^{1/[d(d-1)]}
 +M^{1-2/[d(d-1)]}\lambda^{-2/[d^2(d-1)]}\right).
\]
If $\lambda<M^{-d}$ the last term is at least $M$. Otherwise
$\lambda^{-1/d}\le M$, and the last displayed term bounds it. It can therefore be removed in both cases. Substitute
$\lambda\asymp FM^{-d}$ to obtain \eqref{eq:basederivative}.
For a shorter subinterval first use its own length in this argument;
each term is increasing in that length at fixed $\lambda$, so replacing
it by $M$ is valid. The argument uses the phase on its given interval.
\end{proof}


\section{Theorems used}\label{app:base}

The preceding appendices derive the parts of \cite{B38} used here and
the reductions of de Faveri and Pandey. The analytic input theorems
and the arguments supplied here are identified individually below.

\subsection{Local zeta values and global moments}
Proposition~\ref{prop:localzetaproof} proves the local zeta estimate
\eqref{eq:baselocal}, including the functional-equation and growth
facts used in the proof. Proposition~\ref{prop:momentbaseproof} deduces
both $M_4(T)\prec T$ and $M_{12}(T)\prec T^2$ from the single sparse
fourth-moment input \eqref{eq:basefourth}. This cited theorem gives
the common analytic input for both global moments.

\subsection{Atkinson's short second moment}
With $T^\epsilon\le G\le T^{1/2-\epsilon}$ and
$S_t(\Delta,N)=\sum_{N<n\le N+\Delta}(-1)^nd(n)e(f(t,n))$, the input is
\begin{equation}\label{eq:baseatkinson}
 \begin{split}
 \int_{t-G}^{t+G}|\zeta(1/2+iu)|^2du
 &\ll G\log T\\
 &\quad+G\sum_{\substack{T^{1/3}\le N\ll T G^{-2}(\log T)^{1+\epsilon}\\
                         N\text{ dyadic}}}
 (TN)^{-1/4}\max_{0\le\Delta\le N}|S_t(\Delta,N)|.
 \end{split}
\end{equation}
This is the upper-bound consequence of \cite[Lemma 7.1]{dFP}.
Proposition~\ref{prop:atkinsondirect} proves it directly from the
double Dirichlet series, Euler summation, and Gaussian smoothing.
The estimate uses the exact phase at the center $t$ and includes all
Fourier indices and both nonstationary endpoints. For $t\asymp T$, changing
between the two comparable scales in this display costs only fixed
constants in the dyadic range.

\subsection{Sparse fourth moments}
For disjoint intervals $I_1,\ldots,I_J\subset[T,2T]$ of common length
$\sqrt T\le L\le T$, the input used in the pair count is
\begin{equation}\label{eq:basefourth}
 \sum_{j\le J}\int_{I_j}|\zeta(1/2+it)|^4dt
 \prec JL+T\sqrt{J/L}.
\end{equation}
It is Iwaniec's theorem \cite[Theorem~4]{Iwa80}, as stated in \cite[Proposition 8.1]{dFP}.
Proposition~\ref{prop:pairproof} derives the pair-count consequence
of this cited theorem.

\subsection{Heath-Brown's pair and Vinogradov mean values}
The pair in Lemma~\ref{lem:hb} is Theorem 2 of \cite{HB17}, following
there from Theorems 1 and 3. The underlying derivative estimate is
\begin{equation}\label{eq:basederivative}
 \left|\sum_{m\in I}e(g(m))\right|
 \ll_{d,\epsilon}M^{1+\epsilon}
 \left(F^{1/[d(d-1)]}M^{-1/(d-1)}
       +M^{-1/[d(d-1)]}+F^{-2/[d^2(d-1)]}\right),
\end{equation}
when $|g^{(d)}|\asymp FM^{-d}$ on an interval $I\subset(M,2M]$.
Proposition~\ref{prop:hbderivativeproof} proves this estimate from the
precise optimal Vinogradov mean-value input \eqref{eq:vmvtinput}.
The coefficient-box count, localization, all powers of the shift length
and the endpoint terms are included in that proof. The mean-value
input is the optimal Vinogradov mean-value theorem \cite{BDG16,Woo16}.

The finite-derivative conditions can then be settled explicitly
for the pair needed in this paper. For $F\ge M^{85/8}$ the target
$M^{1-a}F^{1/245}$ is at least $M$, so the trivial bound suffices.
For $M\le F\le M^2$, Weyl differencing and
Lemma~\ref{lem:secondproof}, with differencing length
$\asymp M F^{-1/3}$, give
$O(M^{1/2}F^{1/6}+M F^{-1/6})$; both terms are smaller than the
target. The differencing inequality follows by averaging the $H$
translates of the sum and applying Cauchy; in its usual form it is
\[
 |S|^2\ll M^2/H+(M/H)\sum_{h\le H}
 \left|\sum_m e(g(m+h)-g(m))\right|.
\]
For the reciprocal perturbations here the differenced second derivative
has size $hF/M^3$, so substitution proves the stated bound, including
end intervals. Constants in $H$ can be fixed to keep $m+h$ on a fixed
enlargement.

For $2\le\tau=\log F/\log M\le85/8$, apply
\eqref{eq:basederivative} with the integer $d\in\{3,\ldots,12\}$
whose interval
\[
 d-2+2/d\le\tau\le d-1+2/(d+1)
\]
contains $\tau$. The intervals are consecutive.
On it the maximum of the three relative exponents in
\eqref{eq:basederivative} is at most the chord joining
$-1/[d(d-1)]$ and $-1/[d(d+1)]$ at its two endpoints.
Its slope is $2/((d-1)^2(d+2))$. These slopes decrease with $d$,
so the resulting continuous piecewise linear function is concave and
lies below the line of its $d=8$ segment. That line is
\[
 \frac{\tau}{245}+\frac{269}{280}-\frac1{245}-1.
\]
This proves the desired exponent-pair estimate from
\eqref{eq:basederivative}, using derivatives only through order 12.
The exact phase estimates in Proposition~\ref{prop:exactphaseproof}
and Proposition~\ref{prop:tuplemorseproof} give them with relative
errors $O_d(|\eta|)$ through the specified finite derivative order.

The sixth-moment proof uses \eqref{eq:basefourth} and
\eqref{eq:vmvtinput}. Mean values of degree at most 11 cover the
chosen exponent pair.

The additional estimates in \cite{Supp} use direct differencing for
an alternative eighth-moment bound. The sixth-moment proof here uses
the two stated input theorems.

\section{The needed Atkinson bound by Gaussian smoothing}
\label{app:atkinsondirect}

This appendix proves \eqref{eq:baseatkinson} by Gaussian smoothing.
We start from Atkinson's dissection of the double Dirichlet series
\cite{Atk49,SS23}. Gaussian integration is performed before estimating
the resulting Fourier integrals. This avoids the divisor-error integrals
needed to recover the full unsmoothed formula. The zeta functional
equation and elementary gamma facts used in Appendix~\ref{app:localzeta}
remain available; the dissection and its continuation are supplied here.

\begin{proposition}[short second moment]
\label{prop:atkinsondirect}
Let $\epsilon>0$, $T\ge2$, and
$T^\epsilon\le G\le T^{1/2-\epsilon}$. With the exact phase
\[
 f(T,n)=\frac T\pi\left(\operatorname{arsinh}\sqrt{\frac{\pi n}{2T}}
       +\sqrt{\frac{\pi n}{2T}\left(1+\frac{\pi n}{2T}\right)}\right)
\]
and $S_T(\Delta,N)=\sum_{N<n\le N+\Delta}(-1)^nd(n)e(f(T,n))$,
\[
 \int_{T-G}^{T+G}|\zeta(1/2+it)|^2dt
 \ll_\epsilon G\log T+
 G\sum_{\substack{N\text{ dyadic}\\T^{1/3}\le N\le
 T G^{-2}(\log T)^{1+\epsilon}}}
 (TN)^{-1/4}\max_{0\le\Delta\le N}|S_T(\Delta,N)|.
\]
Fixed endpoint conventions and bounded initial values of $T$ change
only the implied constant.
\end{proposition}

\subsection{Continuation of the dissection}
Put $q=u+v$. Initially for $\Re u,\Re v>1$,
\[
 \zeta(u)\zeta(v)=\zeta(q)+F(u,v)+F(v,u),\qquad
 F(u,v)=\sum_{m,n\ge1}n^{-u}(n+m)^{-v}.
\]
For $\Re v>-1$, define the continued tail by
\[
 R_v(n)=\zeta(v,n+1)-\frac{n^{1-v}}{v-1}+\frac12n^{-v},
\]
where the expression at $v=1$ is interpreted after removing its pole.
Integrate $x^{-v}$ over each
unit interval twice by parts. The elementary periodic Bernoulli
polynomial $B_2(\{x\})=\{x\}^2-\{x\}+1/6$ gives
\begin{equation}\label{eq:EulerR}
 R_v(n)=\frac v{12}n^{-v-1}
   -\frac{v(v+1)}2\int_n^\infty B_2(\{x\})x^{-v-2}\,dx.
\end{equation}
One can first sum to an integer $K$ with $\Re v>1$ and then let
$K\to\infty$; the displayed integral continues the result to
$\Re v>-1$. In particular, on closed substrips,
\[
 |R_v(n)|\ll (1+|v|)^2n^{-\Re v-1}.
\]
It follows by locally uniform convergence that
\begin{equation}\label{eq:Fcontinue}
 F(u,v)=\frac{\zeta(q-1)}{v-1}-\frac12\zeta(q)
                              +\sum_{n\ge1}n^{-u}R_v(n)
\end{equation}
is meromorphic for $\Re q>0$, $\Re v>-1$. This continues the
dissection in the region also satisfying $\Re u>-1$.
The gamma integral gives
\[
 \int_0^\infty y^{-u}(1+y)^{-v}dy
       =\frac{\Gamma(1-u)\Gamma(q-1)}{\Gamma(v)}
\]
when both integrals converge. This follows directly from the double
integral defining the two gamma factors, with variables $x+y$ and
$x/(x+y)$; it then continues meromorphically.
Set
\[
 G(u,v)=F(u,v)-\frac{\Gamma(1-u)\Gamma(q-1)}{\Gamma(v)}\zeta(q-1).
\]
Let $v=1-u+\delta$ and let $\delta\to0$.
Use $\zeta(1+\delta)=\delta^{-1}+\gamma+O(\delta)$,
$\zeta(\delta)=-1/2-\tfrac12\log(2\pi)\delta+O(\delta^2)$,
and $\Gamma(\delta)=\delta^{-1}-\gamma+O(\delta)$.
The zeta constants follow from Euler summation and the functional
equation proved in Appendix~\ref{app:localzeta}; expanding its gamma
factor at zero gives $\zeta'(0)=-\log(2\pi)/2$.
The poles in $\delta$ cancel. Writing $\psi=\Gamma'/\Gamma$ gives
\begin{equation}\label{eq:Gexplicit}
 \mathcal G(u):=G(u,1-u)
 =\frac1{2u}-\gamma+\frac12\log(2\pi)-\frac12\psi(1-u)
                      +\sum_{n\ge1}n^{-u}R_{1-u}(n).
\end{equation}
The last series converges locally uniformly for $\Re u<2$ and is
$O((1+|\Im u|)^2)$ on each closed vertical strip to the left of $1$.
Thus $\mathcal G$ is meromorphic on $\Re u<1$, with sole pole
$u=0$, of residue $1/2$. Combining the two gamma quotients in the
dissection yields, for $0<\Re u<1$,
\begin{equation}\label{eq:exactdissection}
 \zeta(u)\zeta(1-u)=
 \tfrac12\{\psi(u)+\psi(1-u)\}+2\gamma-\log(2\pi)
                         +\mathcal G(u)+\mathcal G(1-u).
\end{equation}
Indeed the gamma terms have constant coefficient
$\tfrac12(\psi(u)+\psi(1-u))+\gamma-\log(2\pi)$ and
pole $-1/\delta$; the term $\zeta(1+\delta)$ adds $1/\delta+\gamma$.
This also checks the normalization of the main term.

\subsection{A Fourier representation on the left line}
For $-1<\Re u<0$ and initially $\Re q>2$, apply Poisson summation
to the even extension of $y^{-u}(1+y)^{-v}$ on the positive half-line,
at mesh $1/m$. Its value at zero is zero. The Fourier series is
absolutely convergent: a split at reciprocal frequency followed by
two integrations by parts gives decay $O(k^{\Re u-1})$.
The periodized function is continuous, so its absolutely convergent
Fourier series evaluates to it. Summing over $m$ gives
\[
 G(u,v)=2\sum_{m,k\ge1}m^{1-q}
       \int_0^\infty y^{-u}(1+y)^{-v}\cos(2\pi kmy)\,dy.
\]
Grouping $n=km$ replaces the coefficient by
$\sum_{m\mid n}m^{1-q}$. This identity continues to
$-1<\Re u<0$, $\Re v>1$, $\Re q>0$: at zero the integral has
Fourier size $O(n^{\Re u-1})$, and at infinity two integrations by
parts contribute $O(n^{-2})$. The divisor coefficient is
$\ll_\delta n^{\max(0,1-\Re q)+\delta}$, so the grouped series
converges locally uniformly exactly on closed smaller regions of
these inequalities. Bounds for the coefficients follow directly from
the divisor sum; the elementary $d(n)\ll_\delta n^\delta$ follows
by separating the finitely many primes below $2^{1/\delta}$ in the
prime factorization. Analytic continuation therefore permits $q=1$.
We obtain
\begin{equation}\label{eq:leftG}
 \mathcal G(u)=\sum_{n\ge1}d(n)H(u,n),\qquad
 H(u,n)=2\int_0^\infty y^{-u}(1+y)^{u-1}\cos(2\pi ny)\,dy.
\end{equation}
At infinity the integral is an ordinary improper oscillatory integral.
For $\Re u=-a$, $0<a<1$, the same split at $1/n$ and two integrations
by parts give
$|H(u,n)|\ll_a(1+|u|)^2 n^{-1-a}$.
This bound justifies integration against a Gaussian and summation
in $n$ on that vertical line. The Fourier bound used
above can be seen in the same way: on $[0,1/n]$ integrate the
absolute value; on $[1/n,1]$ integrate twice, bounding the second
derivative by a constant times $(1+|u|+|v|)^2y^{-\Re u-2}$;
on $[1,\infty)$ the differentiated amplitude is integrable.
All boundary terms have the same stated sizes.

\subsection{An exact smoothed identity}
Let $w_T(t)=\exp(-(t-T)^2/G^2)$ and
$W_T(u)=\exp((u-1/2-iT)^2/G^2)$.
In the integral of $\mathcal G(1/2+it)$ shift the vertical line from
$1/2$ to $-a$. Formula \eqref{eq:Gexplicit} gives polynomial growth;
the horizontal sides tend to zero by Gaussian decay. The sole crossed
residue contributes $\pi W_T(0)$. On the left line use
\eqref{eq:leftG}, and evaluate the elementary Gaussian Fourier integral:
\[
 \frac1i\int_{\Re u=-a}W_T(u)e^{uL}\,du
   =\sqrt\pi G\exp(L/2+iTL-G^2L^2/4),\qquad L>0.
\]
For example substitute $u=-a+it$, complete the square, and shift the
entire Gaussian contour; its horizontal sides again vanish.
In the other term of \eqref{eq:exactdissection} substitute $u\mapsto1-u$;
this gives the conjugate expression with center $-T$.
Thus, with $L(y)=\log((1+y)/y)$,
\begin{equation}\label{eq:exactGaussian}
\begin{split}
 \int_{\mathbb R}|\zeta(1/2+it)|^2w_T(t)dt
 &=\int_{\mathbb R}
   \{\Re\psi(1/2+it)+2\gamma-\log(2\pi)\}w_T(t)dt\\
 &\quad+4\sqrt\pi G\Re\sum_{n\ge1}d(n)\mathcal I_n(T,G)
       +2\pi\Re W_T(0),\\
 \mathcal I_n(T,G)&=\int_0^\infty
 \frac{e^{-G^2L(y)^2/4}}{\sqrt{y(1+y)}}e^{iTL(y)}
                                    \cos(2\pi ny)\,dy.
\end{split}
\end{equation}
The exchanges follow from the following explicit tail bounds:
first truncate the $y$ integral, perform the Gaussian integral, and
then use one integration by parts on its tail. The tail before
Gaussian integration is $O((1+|u|)/(nY))$ for $y\ge Y\ge1$,
which is integrable in $u$ against the Gaussian. Summation and the
$u$ integral were interchanged using the absolutely summable
bound following \eqref{eq:leftG}.
After Gaussian integration the Fourier amplitude is flat at $y=0$,
and its differentiated amplitudes are integrable at infinity. Two
integrations by parts show absolute summability of the final $n$ series
for each fixed $T,G$. This also verifies the final identity directly.

The first integral in \eqref{eq:exactGaussian} is $O(G\log T)$:
$\Re\psi(1/2+it)=O(\log(2+|t|))$ by the gamma estimate
proved in Appendix~\ref{app:localzeta}. Also
$|W_T(0)|=\exp((1/4-T^2)/G^2)$ is negligible.
It remains to estimate the kernels $\mathcal I_n$, keeping their phases.

\subsection{Stationary and nonstationary terms}
For $1\le n\le T$ let $a>0$ be defined by
\[
 \frac{a^2}{1+a}=\frac{2\pi n}{T},\qquad Y=Ta,\qquad b=Ga.
\]
Here $a$ belongs to a fixed bounded interval,
$Y\asymp\sqrt{nT}$, and $b^2\asymp G^2n/T$.
Put $y=x/a$ and split the cosine into its two exponentials.
Their phases divided by $Y$ and common amplitude are
\[
 \varphi_\pm(x,a)=a^{-1}\log(1+a/x)\pm x/(1+a),\qquad
 A(x,a,b)=\frac{\exp[-b^2\{a^{-1}\log(1+a/x)\}^2/4]}
                    {\sqrt{x(x+a)}}.
\]
These expressions extend smoothly to $a=0$ for $x$ bounded away
from zero. The plus phase has its unique critical point at $x=1$,
with $\varphi_+''(1,a)=(2+a)/(1+a)^2$; the minus phase has none.
Choose a fixed smooth cutoff $\chi$ supported in $[1/2,2]$ and equal
to one on $[3/4,3/2]$.

Apply the exact Morse--Fresnel identity of Lemma~\ref{lem:exactmorse}
to $\chi A e^{iY\varphi_+}$. Replace the scale $Y$ by $Y/(2\pi)$
to pass from the lemma's $e$ notation to the radian phase. It gives an exact, complex coefficient $c(T,G,n)$ with
\begin{equation}\label{eq:stationaryAtkinson}
 \int\chi(x)A(x,a,b)e^{iY\varphi_+(x,a)}dx
    =e^{i(2\pi f(T,n)-\pi n)}c(T,G,n),
\end{equation}
where
\begin{equation}\label{eq:Atkinsonsymbol}
 |c|+n|\partial_nc|\ll(Tn)^{-1/4},\qquad
 |c|\ll(Tn)^{-1/4}e^{-c_1G^2n/T}.
\end{equation}
Here and below $c_1>0$ denotes a fixed decay constant.
To verify all uniform bounds, on the support of $\chi$ the function
$a^{-1}\log(1+a/x)$ is bounded above and below by positive constants.
Every fixed $x$, $a\partial_a$, and $b\partial_b$ derivative of the
amplitude is bounded by a polynomial in $b$ times $e^{-c_2b^2}$,
hence by $C_j e^{-c_3b^2}$. The Morse substitution and its inverse
have uniformly bounded derivatives on this fixed compact set.
In the Fourier form of the exact Fresnel identity, both parameter
derivatives and $Y\partial_Y$ are bounded by finitely many Fourier
moments of this amplitude, proving the same bounds for $Y^{1/2}c$.
Finally
$n\partial_na/a=(1+a)/(2+a)$, and the logarithmic derivatives of
$Y$ and $b$ are the same bounded number. This proves
\eqref{eq:Atkinsonsymbol}, including its variation bound. The phase
at $x=1$ is
$T\log(1+a)+Ta/(1+a)=2\pi f(T,n)-\pi n$, exactly.

The omitted parts cost $O(Y^{-1})$, uniformly in $G$. Here is a
first-derivative proof that also handles both improper endpoints.
Write $t_x=x(x+a)$. For the lower nonstationary part of the plus
phase the absolute amplitude divided by the phase derivative is
\[
 \frac1Y\frac{\sqrt{t_x}}{1-t_x/(1+a)}
           \exp[-G^2\log(1+a/x)^2/4],\qquad x\le3/4.
\]
All three positive factors increase with $x$, and their value at
$3/4$ is $O(Y^{-1})$. Its total variation is therefore $O(Y^{-1})$.
For $x\ge3/2$ the rational factor becomes
$\sqrt{t_x}/(t_x/(1+a)-1)$, which decreases to zero, while the
Gaussian factor increases to one. The variation of their product
is at most the sum of the two variations, again $O(Y^{-1})$.
The cutoff transition intervals have the same bound. For the minus
phase the rational factor is $\sqrt{t_x}/(1+t_x/(1+a))$; it is
unimodal, vanishes at both ends, and has bounded maximum.
Its product with the increasing Gaussian has bounded variation.
One integration by parts proves the claimed $O(Y^{-1})$ bounds in
all cases; endpoint terms vanish. Hence
\begin{equation}\label{eq:kernelAtkinson}
 \mathcal I_n(T,G)=\tfrac12(-1)^ne(f(T,n))c(T,G,n)
                           +O((nT)^{-1/2}),\qquad 1\le n\le T.
\end{equation}
The sum of the errors against $d(n)$ is $O(\log T)$, since the
divisor identity $d(n)=\sum_{ab=n}1$ gives
$\sum_{n\le X}d(n)n^{-1/2}\ll\sqrt X\log(2X)$.

All $n>T$ cost $O(\log T/T)$.
Take a fixed smooth partition at $y=1$, with its lower piece supported
in $y\le2$. On this piece the amplitude and its first two derivatives,
including those of $e^{iTL(y)}$, have integrable absolute values
$\ll (1+T)^2(1+G)^C e^{-c_4G^2}$.
This follows near zero by substituting $s=\log(1/y)$: every resulting
power of $1/y$ becomes an exponential in $s$, dominated by
$e^{-G^2s^2/4}$; on $[1/2,2]$ use $L(y)\ge\log(3/2)$.
Two integrations by parts against $\cos(2\pi ny)$ bound this piece
by $O_\epsilon(n^{-2})$ for $G\ge T^\epsilon$.
On the upper piece $y\ge1$, the phases
$TL(y)\pm2\pi ny$ have first derivative of modulus $\gg n$
and higher derivatives $O_j(Ty^{-j-1})$.
For $j=1,2$ the common amplitude has derivative
$O_j(y^{-1-j})$, uniformly in $G$: each Gaussian derivative gives
powers of $G/y$ times $e^{-c(G/y)^2}$, which are bounded.
Two integrations by parts therefore cost
$O(n^{-2}(1+T/n+(T/n)^2))=O(n^{-2})$.
After the first integration all required derivatives are integrable and the
boundary term tends to zero. Finally
$\sum_{n>T}d(n)n^{-2}\ll\log T/T$, again by the divisor identity.

\subsection{Truncation and the partial sums}
Substitute \eqref{eq:kernelAtkinson} in \eqref{eq:exactGaussian}.
The nonstationary terms, the diagonal gamma integral, and the residue
are all within $O_\epsilon(G\log T)$.
For $n>TG^{-2}(\log T)^{1+\epsilon}$ the exponential estimate in
\eqref{eq:Atkinsonsymbol} makes the stationary contribution smaller
than any fixed negative power of $T$ after summing $n\le T$.
For $n\le2T^{1/3}$ the trivial divisor bound in mean gives
\[
 \sum_{n\le2T^{1/3}}d(n)(Tn)^{-1/4}\ll\log T.
\]
On each remaining dyadic interval, Abel summation and
\eqref{eq:Atkinsonsymbol} give
\[
 \left|\sum_{N<n\le N'}(-1)^nd(n)e(f(T,n))c(T,G,n)\right|
 \ll (TN)^{-1/4}\max_{0\le\Delta\le N}|S_T(\Delta,N)|,
 \qquad N'\le2N.
\]
All coefficients and the phase are those of the sums $S_T(\Delta,N)$.
The center remains the height $T$ throughout. Finally $w_T(t)\ge e^{-1}$ for $|t-T|\le G$.
Positivity bounds the requested short integral by $e$ times the full
Gaussian integral. This proves Proposition~\ref{prop:atkinsondirect}
and hence \eqref{eq:baseatkinson}.

\begin{remark}
The proof establishes the Gaussian-smoothed upper-bound consequence
used in this paper. Attribution for the
dissection and the short-second-moment approach remains with
Atkinson~\cite{Atk49}, Heath-Brown~\cite{HB78} and Ivi\'c. The advantage here is that smoothing
precedes the Fourier-integral estimates, allowing the needed input
to follow from the displayed dissection and kernel estimates.
\end{remark}

\section{Computations}\label{app:check}
The file \path{checks/exponents.py} checks the nine exponents,
their intersection, the parameter ranges, and the sixth-moment sum.
The file \path{checks/identity_checks.py} checks the phase identities.
These calculations use exact arithmetic.

The file \path{checks/atkinson_checks.py} checks the phase formulas
and compares the Gaussian identity with numerical integration at two
values of $T$.

\section*{Use of AI tools}
AI tools were used for proof search, calculations, and writing.

\begin{thebibliography}{99}

\bibitem[Atk49]{Atk49}
F.~V. Atkinson, \emph{The mean-value of the Riemann zeta function}, Acta Math. \textbf{81} (1949), 353--376.

\bibitem[SS23]{SS23}
A.~Simoni\v{c} and V.~V. Starichkova,
\emph{Atkinson's formula for the mean square of $\zeta(s)$ with an explicit error term},
arXiv:2105.06821v3 (2022), published 2023.

\bibitem[B38]{B38}
bramble\_trace (AI agent), \emph{Curved Atkinson phases, reciprocal cancellation and the sixth moment of the Riemann zeta function}, rStation paper~38, version~3, 21 October 2026, \url{https://rstation.ai/problems/riemann_hypothesis/papers/38}.

\bibitem[BDG16]{BDG16}
J.~Bourgain, C.~Demeter and L.~Guth, \emph{Proof of the main conjecture in Vinogradov's mean value theorem for degrees higher than three}, Ann. of Math. (2) \textbf{184} (2016), 633--682.

\bibitem[CG98]{CG98}
J.~B. Conrey and A.~Ghosh, \emph{A conjecture for the sixth power moment of the Riemann zeta-function}, Internat. Math. Res. Notices (1998), no.~15, 775--780.

\bibitem[dFP]{dFP}
A.~de~Faveri and M.~Pandey, \emph{The sixth moment of the Riemann zeta function}, arXiv:2610.01035v1 (2026).

\bibitem[HB78]{HB78}
D.~R. Heath-Brown, \emph{The twelfth power moment of the Riemann zeta-function}, Quart. J. Math. Oxford Ser. (2) \textbf{29} (1978), 443--462.

\bibitem[HB17]{HB17}
D.~R. Heath-Brown, \emph{A new $k$-th derivative estimate for exponential sums via Vinogradov's mean value}, Proc. Steklov Inst. Math. \textbf{296} (2017), 88--103; arXiv:1601.04493.

\bibitem[Ing26]{Ing26}
A.~E. Ingham, \emph{Mean-value theorems in the theory of the Riemann zeta-function}, Proc. London Math. Soc. (2) \textbf{27} (1926), 273--300.

\bibitem[Iwa80]{Iwa80}
H.~Iwaniec, \emph{Fourier coefficients of cusp forms and the Riemann zeta-function}, S\'eminaire de Th\'eorie des Nombres de Bordeaux (1979--1980), Exp.~No.~18, 36~pp.

\bibitem[TY]{TY}
T.~S. Trudgian and A.~Yang, \emph{Toward optimal exponent pairs}, Math. Comp. \textbf{94} (2025), 1467--1502, \url{https://doi.org/10.1090/mcom/3993}.

\bibitem[Woo16]{Woo16}
T.~D. Wooley, \emph{The cubic case of the main conjecture in Vinogradov's mean value theorem}, Adv. Math. \textbf{294} (2016), 532--561.

\bibitem[Thin]{Thin} M.~Chae,
\emph{Thin planar energy and the eighth moment of the Riemann zeta function},
17 November 2026, \href{https://github.com/chaemin333-max/riemann-zeta-eighth-moment/blob/main/paper/moments.pdf}{PDF}.

\bibitem[Supp]{Supp} M.~Chae,
\emph{Additional estimates for zeta moments},
17 November 2026, \href{https://github.com/chaemin333-max/riemann-zeta-eighth-moment/blob/main/supplement/moments.pdf}{PDF}.

\end{thebibliography}

\end{document}

